
% BEGIN preamble.tex
\documentclass[11pt,letterpaper]{article}
\usepackage[T1]{fontenc}
\usepackage[utf8]{inputenc}
\usepackage[sc]{mathpazo}
\usepackage[margin=0.93in,headheight=15pt]{geometry}
\usepackage{amsmath,amssymb,amsthm,mathtools}
\usepackage{microtype}
\usepackage{booktabs,longtable,array}
\usepackage{enumitem}
\usepackage{xcolor}
\usepackage{xurl}
\usepackage{fancyhdr}
\usepackage{hyperref}
\hypersetup{colorlinks=true,linkcolor=blue!45!black,citecolor=green!35!black,
 urlcolor=blue!50!black,
 pdftitle={Boundary Contacts and Logarithmic Spectral Integrals},
 pdfauthor={Research continuation prepared for Vladimir Reshetnikov with OpenAI assistance},
 pdfsubject={Independent Hurwitz-order contacts, all-mass spectral jets, Gamma, polygamma, harmonic and polylogarithmic identities}}
\numberwithin{equation}{section}
\newtheorem{theorem}{Theorem}[section]
\newtheorem{lemma}[theorem]{Lemma}
\newtheorem{proposition}[theorem]{Proposition}
\newtheorem{corollary}[theorem]{Corollary}
\theoremstyle{definition}
\newtheorem{definition}[theorem]{Definition}
\newtheorem{question}[theorem]{Research question}
\theoremstyle{remark}
\newtheorem{remark}[theorem]{Remark}
\DeclareMathOperator{\Li}{Li}
\DeclareMathOperator{\Log}{Log}
\DeclareMathOperator{\sech}{sech}
\DeclareMathOperator{\FP}{FP}
\newcommand{\dd}{\,d}
\newcommand{\src}[1]{\nolinkurl{#1}}
\newcommand{\pin}{\texttt{e3d7ec5d3f9be259c1f91edaea4389f4e79b5b03}}
\newcommand{\stir}[2]{\genfrac{[}{]}{0pt}{}{#1}{#2}}
\newcommand{\one}[1]{\{1\}^{#1}}
\allowdisplaybreaks[2]
\setlist{topsep=4pt,itemsep=4pt,leftmargin=*}
\setlength{\emergencystretch}{1.5em}
\setlength{\parskip}{3pt}
\pagestyle{fancy}
\fancyhf{}
\fancyhead[L]{\small Boundary Contacts and Logarithmic Spectral Integrals}
\fancyhead[R]{\small ProveIt research continuation}
\fancyfoot[C]{\thepage}
\renewcommand{\headrulewidth}{0.3pt}
\setcounter{tocdepth}{2}
\title{\textbf{Boundary Contacts and\\Logarithmic Spectral Integrals}\\[7pt]
 \large Independent Hurwitz-order jets, all-mass resolvent formulas,\\
 and harmonic--polylogarithmic consequences}
\author{Research continuation prepared for Vladimir Reshetnikov\\
 \small ProveIt / Analysis / Polylogarithms\\
 \small Mathematical derivations and computational checks with OpenAI assistance}
\date{11 October 2026}

% END preamble.tex

\begin{document}
\maketitle
\begin{abstract}
We continue two explicit lines of research in the ProveIt manuscript and
its incoming reports. First, we complete the contact algebra for the
Gamma-normalized product of two independently ordered and independently
shifted alternating Hurwitz tails. At every hyperplane \(p+q=-2m\), the
symmetric cutoff contact and the alternating parity contact are finite
rational polynomials of total spectral degree at most \(m\). We give a
terminating coefficient algorithm, every independent mixed order
derivative, exact cutoff translation, a sharp highest-degree generating
function, and ordinary convergent polynomial or logarithmic subtractions.
Two concrete series on the first negative-total-order hyperplane evaluate
in \(\pi\), \(\log 2\), and rational numbers.

Second, every derivative at zero of a scalar spectral zeta function with
any number of quadratic masses is expressed by one absolutely convergent
real integral. Its density is an explicit polynomial in real logarithms
and \(\pi\), and its kernel is a digamma difference. Finite asymptotic subtraction
extends the theorem to every even polynomial weight and, consequently,
to every spectral derivative at every nonpositive integer. Products of logarithmic
mass increments have a compactly supported density with exact vanishing
moments. This gives all connected interactions, regular coincident-mass
limits, a critical odd-weight subtraction, polygamma derivatives, and
anchored log-Gamma primitives. The repeated-mass density has a beta
generating function; its moments are Stirling and generalized harmonic
polynomials, and its polylogarithmic transform is a height-one multiple
polylogarithm.

The proofs specify branches, continuation, cutoff coordinates, and every
interchange. We retain the attribution of the existing local spectral
correction theorem and of classical Gamma, Hurwitz, and multiple-zeta
identities. A source audit supplies a definite reciprocal-Gamma
transcription correction and reconciles two stale research targets with
already-preserved proofs. Exact symbolic checks and separate high-precision
diagnostics accompany the article. Historical priority, arithmetic
independence, interval certification, and proof-assistant formalization are
not inferred from those checks.
\end{abstract}
\noindent\textbf{Keywords.}
Alternating Hurwitz zeta; cutoff finite part; spectral zeta derivative;
Gamma and polygamma; generalized Stieltjes constants; generalized harmonic
numbers; logarithmic density; multiple polylogarithm.
\clearpage
\begingroup
\small
\setlength{\parskip}{0pt}
\tableofcontents
\endgroup
\clearpage

% BEGIN sections/01_scope.tex
\section{Scope, main advances, and source status}
\label{bc:scope}

This article is a continuation of the ProveIt material at the immutable
repository checkpoint
\[
 \pin .
\]
The requested manuscript directory, the seven archives in
\src{docs/incoming}, and the directly relevant preserved continuations were
inspected together. This matters: an open question in one report may have
been solved in another report that has not yet been integrated into the
canonical manuscript. The accompanying provenance file records the
identities and hashes of the sources used here. The source tree itself was
not modified.

\subsection{The questions answered}

The first precise target is the question headed \emph{Every bilinear contact
hyperplane} in \emph{Bilinear Tails and Determinants}
\cite{ProveItBilinear}. Its finite two-order endpoint identity is already
proved. The missing data are the bivariate contact polynomials on
\(p+q=-2m\), all independent order derivatives, and the changes under
translation of the cutoff. A diagonal restriction cannot recover
derivatives transverse to that restriction. Section~\ref{ct:section}
supplies these data for every \(m\ge0\), including explicit ordinary
convergent representatives. The finite endpoint law and its first
\(1/2\) contact remain attributed to the predecessor.

The second target is the higher-order, many-mass question in
\emph{Higher Quadratic Jets}, item 3, and \emph{Beta Polarization of Higher
Spectral Defects}, item 3
\cite{ProveItHigherQuadratic,ProveItHigherDefects}. The latter already gives
a complete pairwise formula for second jets. Its research question asks
for higher connected logarithmic terms and correctly normalized
coincident-mass limits. Section~\ref{sp:section} gives a single real
resolvent integral for every scalar jet and a compact integral for every
connected logarithmic product. The number of real integration variables
stays one as both the number of masses and the derivative order increase.
Section~\ref{ng:section} further extends the full-jet formula to every
even polynomial weight and to every nonpositive integer spectral point.
This settles the scalar representation and connected-product parts of the
question, and supplies an alternative explicit representation for the
higher negative-point jets. It does not assert that all those integrals reduce to finitely
many special values of a smaller function class.

\begin{center}
\small
\renewcommand{\arraystretch}{1.25}
\begin{tabular}{>{\raggedright\arraybackslash}p{0.22\textwidth}>{\raggedright\arraybackslash}p{0.31\textwidth}>{\raggedright\arraybackslash}p{0.38\textwidth}}
\toprule
Topic & Result in this article & Status relative to inspected sources\\
\midrule
Two-order tail contacts &
Finite coefficients at every resonance; all mixed jets and translations &
Completes the explicitly open bivariate contact question.\\
Ordinary boundary sums &
Polynomial and logarithmic subtractions, with exact examples &
Consequences of the completed contact algebra.\\
Full scalar spectral jets &
One digamma integral for every \(N,r\) &
Explicit all-order representation beyond the supplied second-jet formulas.\\
Even weights and negative points &
Every polynomial degree and every derivative at each nonpositive integer &
A finite Hurwitz baseline and the same single ordinary integral.\\
Connected logarithmic products &
Compact density, moments, critical odd weight, collisions &
All-order realization of the higher connected-term target.\\
Local spectral correction &
Finite residue polynomial at fixed \(N\) &
Inherited mechanism, restated with matching normalization.\\
Harmonic and polylogarithmic forms &
Beta moments, confluent expansions, height-one transform &
Derived bridge to classical coefficient and function alphabets.\\
\bottomrule
\end{tabular}
\end{center}

\subsection{Two organizing mechanisms}

The tail calculation is controlled by an \emph{upper endpoint}. Its
asymptotic coefficient at a resonant exponent gives the discrepancy between
an analytic continuation and a specified cutoff finite part. The necessary
operation is to expand in independent parameters before restricting to the
resonance. A finite-degree theorem then makes all sufficiently high
independent contact derivatives vanish. The same asymptotic coefficients
specify actual convergent subtractions, so the finite-part formulas have
ordinary-series counterparts.

The spectral calculation is controlled by a \emph{logarithmic jump across
a cut}. Subtracting the mass-zero logarithm removes the obstruction at
infinity. Cauchy's formula converts the remaining logarithmic expression
into a resolvent integral, and the lattice resolvent sums exactly to a
digamma difference. For a product of logarithmic mass increments the cut
is compact and the decay at infinity forces its first moments to vanish.
Those moment cancellations supply the polynomial subtractions for weighted
sums.

These are different mechanisms. The cutoff degree theorem is a statement
about asymptotic coefficients in the orders \(p,q\). The many-mass
interaction hierarchy is a statement about finite differences of spectral
jets with a fixed total number of factors. Neither follows by replacing
one meaning of degree with the other.

\subsection{Classical inputs and the priority boundary}

The Gamma quotient expansion, the alternating Hurwitz expansion and its
parameter derivatives, the digamma resolvent identity, and Euler's beta
integral are classical
\cite{DLMF,HuKim,WangHuKim}. Products of zeta tails have an established
multiple-zeta literature; Hoffman's paper is a relevant example
\cite{Hoffman}. The broad relation between regularized traces and local
asymptotic coefficients is likewise established
\cite{PaychaScott}. The normalized exponent-local degree theorem is already
present in the project's \emph{Arithmetic Spectral Jets} continuation
\cite{ProveItArithmeticJets}. We use it for comparison and give a direct
normalization check where it enters.

The contribution claimed here is the stated collection of precise
deductions and completed research targets relative to the inspected
project material. The article gives proofs of those statements. A targeted
search did not establish an earlier occurrence of every displayed form,
but that is not a comprehensive historical-priority result. In particular,
we do not relabel a classical beta integral, Stirling-number identity, or
height-one multiple-zeta reduction as a newly discovered constant.

\subsection{Conventions}

\(\zeta^{(r)}(s,a)\) always means an \emph{order} derivative
\(\partial_s^r\zeta(s,a)\). Superscripts on \(\psi\) mean argument
derivatives. The generalized Stieltjes convention is
\[
 \zeta(1+u,a)=\frac1u+
   \sum_{j\ge0}\frac{(-1)^j}{j!}\gamma_j(a)u^j,
 \qquad \gamma_0(a)=-\psi(a).
\]
The generalized harmonic numbers are
\[
 H_n^{(j)}=\sum_{k=1}^{n}k^{-j},\qquad H_0^{(j)}=0,
 \qquad H_n=H_n^{(1)}.
\]
A jet of order \(r\) is the actual derivative, with its factorial included;
\([u^r]\) instead denotes an ordinary Taylor coefficient.

Every finite part is tied to its stated coordinate or independent
regulator. For the tail problem, the coordinate is \(X=N+c\) and the
retained term is \(X^0(\log X)^0\). For the alternating endpoint a parity
subsequence is also specified. All integrals labeled ordinary are
absolutely convergent improper integrals unless conditional convergence is
explicitly stated. Signed logarithmic densities are allowed; positivity
is not assumed or required.

% END sections/01_scope.tex


% BEGIN sections/02_contacts.tex
\section{Independent orders, independent shifts, and every bilinear contact}
\label{ct:section}

The incoming \emph{Bilinear Tails and Determinants} report establishes
bilinear identities for two independently shifted alternating Hurwitz tails\cite{ProveItBilinear},
but leaves their higher contact hyperplanes explicitly open. Its diagonal
calculation does not determine mixed derivatives transverse to the diagonal.
We resolve that question by a finite polynomial calculus. In addition to
all cutoff finite parts, the calculation produces ordinary convergent
identities on every negative-total-order hyperplane. The first such family
has the particularly simple correction $1/2$ in every summand.

The finite bilinear identities recalled below are inherited. The new
statements are the centered coefficient algebra, the independent contact
and translation formulas, their sharp spectral degree, and the ordinary
identities extracted from them. The classical analytic inputs are the
Gamma quotient expansion and the alternating Hurwitz Euler expansion;
no assertion of global priority for these ingredients is intended.

\subsection{The inherited endpoint law and its normalization}

On $\Re z>0$ put
\begin{align}
 \mathcal B_s(z)&=2^{-s}\left\{\zeta\left(s,\frac z2\right)
              -\zeta\left(s,\frac{z+1}{2}\right)\right\},
                   \label{ct:B}\\
 h(z)&=\sqrt2\,\frac{\Gamma((z+1)/2)}{\Gamma(z/2)},
 \qquad \ell(z)=\log h(z),\qquad
 \tau_s(z)=e^{s\ell(z)}\mathcal B_s(z),\label{ct:tau}\\
 V_s(z)&=\frac{z^s-(z-1)^s}{h(z)^{2s}},\qquad \Re z>1.
                   \label{ct:V}
\end{align}
All powers use the holomorphic logarithms real on the positive axis.
The common Hurwitz pole in \eqref{ct:B} cancels, so this expression is
entire in $s$. The classical recurrences give
\begin{equation}\label{ct:recurrences}
 h(z)h(z+1)=z,\qquad
 \mathcal B_s(z)+\mathcal B_s(z+1)=z^{-s},\qquad
 \tau_s(z-1)-\tau_s(z+1)=V_s(z)\tau_s(z).
\end{equation}
The last identity follows by multiplying the two adjacent Hurwitz
recurrences by the appropriate Gamma powers and subtracting.

Let $\Re a,\Re b>0$ and let $p,q$ be independent complex orders. Define,
for $n\ge-1$,
\begin{align}
 u_n&=\tau_p(n+a+1),&v_n&=\tau_q(n+b+1),\notag\\
 C_n&=u_nv_{n+1}+u_{n+1}v_n,&
 D_n&=u_nv_{n+1}-u_{n+1}v_n,\label{ct:CD}\\
 E_+&=C_{-1},&E_-&=D_{-1}.\notag
\end{align}
For later convenience, the signs are kept outside the definitions
\[
 g_n^+=\{V_p(n+a+1)+V_q(n+b+1)\}u_nv_n,
 \qquad
 g_n^-=\{V_p(n+a+1)-V_q(n+b+1)\}u_nv_n.
\]
The source's exact finite identities are
\begin{equation}\label{ct:finite}
 S_N^+:=\sum_{n=0}^N g_n^+=E_+-C_N,
 \qquad
 S_N^-:=\sum_{n=0}^N(-1)^ng_n^-=E_-+(-1)^ND_N.
\end{equation}
Indeed, \eqref{ct:recurrences} gives
$g_n^+=C_{n-1}-C_n$ and $g_n^-=D_{n-1}+D_n$.
These identities hold at every complex order before taking a limit.
For $\Re(p+q)>0$ both infinite sums converge absolutely to $E_+$ and
$E_-$ respectively. The entire functions $E_\pm$ are their endpoint
continuations. A cutoff finite part need not equal that continuation.

\subsection{A centered finite coefficient algebra}

Use Bernoulli polynomials with $B_1(x)=x-1/2$ and Euler numbers normalized by
\[
 \operatorname{sech}t=\sum_{j\ge0}E_{2j}\frac{t^{2j}}{(2j)!}.
\]
Thus $E_0=1$, $E_2=-1$, and $E_4=5$. Put
\begin{equation}\label{ct:Lj}
 L_j=\frac{2^{2j}B_{2j+1}(1/4)}{j(2j+1)}\quad(j\ge1),
 \qquad L_1=\frac1{16},\quad L_2=-\frac5{128},\quad L_3=\frac{61}{768}.
\end{equation}
Define the formal series
\begin{equation}\label{ct:Tformal}
 T_s(v)=\exp\left(s\sum_{j\ge1}L_jv^j\right)
       \sum_{j\ge0}\frac{E_{2j}(s)_{2j}}{2^{2j}(2j)!}v^j
       =\sum_{j\ge0}t_j(s)v^j.
\end{equation}
Each coefficient is a finite rational polynomial; the infinite formal
series is not asserted to converge. An explicit terminating recursion is
\begin{align}
 e_0(s)&=1,&
 e_n(s)&=\frac{s}{n}\sum_{j=1}^n jL_je_{n-j}(s),\label{ct:recursion}\\
 t_n(s)&=\sum_{j=0}^ne_j(s)
       \frac{E_{2n-2j}(s)_{2n-2j}}{2^{2n-2j}(2n-2j)!}.\notag
\end{align}
In particular,
\begin{equation}\label{ct:first-t}
 t_0=1,\qquad t_1(s)=-\frac{s(2s+1)}{16},\qquad
 t_2(s)=\frac{s(20s^3+108s^2+211s+60)}{1536}.
\end{equation}

\begin{lemma}[Centered Gamma--Hurwitz expansion]\label{ct:centered}
For each $K\ge0$,
\begin{equation}\label{ct:tauexpansion}
 \tau_s(x+\tfrac12)=\frac12x^{-s/2}
       \sum_{j=0}^Kt_j(s)x^{-2j}
       +O\bigl(x^{-\Re s/2-2K-2}\bigr).
\end{equation}
The expansion is locally uniform in the order and in bounded complex
shifts of $x$ as $x$ tends to positive infinity. Fixed spectral derivatives
have the differentiated expansion, with finite powers of $\log x$;
fixed argument derivatives improve the decay by their order.
Furthermore $\deg_s t_j\le2j$.
\end{lemma}
\begin{proof}
The classical Gamma quotient expansion \cite{DLMF}, together with the alternating Hurwitz expansion \cite{HuKim}, gives
\[
 \log h(x+\tfrac12)=\tfrac12\log x+
                 \sum_{j=1}^K L_jx^{-2j}+O(x^{-2K-2}).
\]
Only even inverse powers remain because
$B_{n+1}(3/4)-B_{n+1}(1/4)=0$ for odd $n$; the surviving coefficient
is precisely \eqref{ct:Lj}. For $\Re s>0$ the alternating Hurwitz
integral is
\[
 \mathcal B_s(x+\tfrac12)=\frac1{2\Gamma(s)}
  \int_0^\infty t^{s-1}e^{-xt}\operatorname{sech}(t/2)\,dt.
\]
Expanding the smooth factor at zero gives
\[
 \mathcal B_s(x+\tfrac12)\sim\frac12x^{-s}
        \sum_{j\ge0}\frac{E_{2j}(s)_{2j}}{2^{2j}(2j)!}x^{-2j}.
\]
For an arbitrary compact order set the same finite expansion, with a
uniform remainder, follows by taking sufficiently many terms in the
ordinary Hurwitz Euler--Maclaurin expansions of \eqref{ct:B}. This avoids
using the displayed integral outside its convergence domain. Multiplying
finite expansions with their remainders proves \eqref{ct:tauexpansion}.
Cauchy estimates on slightly larger compact sets justify spectral
 derivatives; the differentiated Hurwitz and Gamma expansions justify
argument derivatives. Finally $\deg e_j\le j$, while
$\deg(s)_{2n-2j}=2n-2j$, proving the asserted degree bound.
\end{proof}

For an integer $k\ge0$ and a shift $v$ define
\begin{equation}\label{ct:bk}
 b_k(s,v)=\sum_{j=0}^{\lfloor k/2\rfloor}
 t_j(s)\binom{-s/2-2j}{k-2j}v^{k-2j}.
\end{equation}
Throughout, generalized binomial coefficients in such formulas mean their
polynomial expressions. They must not be computed by separately
specializing Gamma quotients at poles. We have
\begin{equation}\label{ct:b-properties}
 \deg_s b_k\le k,\qquad b_k(s,-v)=(-1)^kb_k(s,v).
\end{equation}

Choose a \emph{fixed} cutoff parameter $c$, initially independent of
$p,q,a,b$, and write
\begin{equation}\label{ct:fixed-cutoff}
 X=N+c,\quad \lambda=\frac{p+q}{2},\quad
 A=a+\tfrac12-c,\quad B=b+\tfrac12-c.
\end{equation}
The complete ambient coefficient polynomials are
\begin{align}
 c_k(p,q;a,b;c)
 &=\frac14\sum_{j=0}^k\bigl\{
 b_j(p,A)b_{k-j}(q,B+1)+b_j(p,A+1)b_{k-j}(q,B)\bigr\},
                         \label{ct:ck}\\
 d_k(p,q;a,b;c)
 &=\frac14\sum_{j=0}^k\bigl\{
 b_j(p,A)b_{k-j}(q,B+1)-b_j(p,A+1)b_{k-j}(q,B)\bigr\}.
                         \label{ct:dk}
\end{align}
Thus $c_0=1/2$, $d_0=0$, and both polynomials have total degree at most
$k$ in the independent orders $p,q$. Lemma~\ref{ct:centered} implies
\begin{align}
 C_N&=X^{-\lambda}\sum_{k=0}^Kc_kX^{-k}
          +O(X^{-\Re\lambda-K-1}),\label{ct:Cexp}\\
 D_N&=X^{-\lambda}\sum_{k=0}^Kd_kX^{-k}
          +O(X^{-\Re\lambda-K-1}).\label{ct:Dexp}
\end{align}
These expansions inherit the same uniformity and derivative properties.

\subsection{Every independent contact and every mixed order derivative}

Fix $p_0+q_0=-2m$, where $m\ge0$ is an integer. Differentiate at independent
ambient $p,q$ first, then specialize to $(p_0,q_0)$. The symbol
$\operatorname{FP}_X$ below denotes the coefficient of
$X^0(\log X)^0$ in the resulting asymptotic expansion.

\begin{theorem}[Complete independent contact calculus]\label{ct:contacts}
For every $r,s\ge0$,
\begin{align}
 \operatorname{FP}_X
 \left.\partial_p^r\partial_q^sS_N^+\right|_{(p_0,q_0)}
 &=\left.\partial_p^r\partial_q^sE_+\right|_{(p_0,q_0)}
    -\left.\partial_p^r\partial_q^sc_m\right|_{(p_0,q_0)}.
                                          \label{ct:pluscontact}
\end{align}
For the alternating law there are two parity finite parts. On the
subsequence $N\equiv\epsilon\pmod2$, using the same coordinate $X=N+c$,
\begin{align}
 \operatorname{FP}_{X,\epsilon}
 \left.\partial_p^r\partial_q^sS_N^-\right|_{(p_0,q_0)}
 &=\left.\partial_p^r\partial_q^sE_-\right|_{(p_0,q_0)}
   +(-1)^\epsilon\left.\partial_p^r\partial_q^sd_m\right|_{(p_0,q_0)}.
                                          \label{ct:minuscontact}
\end{align}
Every contact derivative of total order $r+s>m$ vanishes. The same formulas
hold after fixed argument derivatives if $c$ is held independent while
these derivatives are taken.
\end{theorem}
\begin{proof}
For any coefficient polynomial $a_k(p,q)$,
\begin{align}
 \partial_p^r\partial_q^s\{a_kX^{-\lambda-k}\}
 =X^{-\lambda-k}\sum_{i=0}^r\sum_{j=0}^s
 \binom ri\binom sj(\partial_p^i\partial_q^ja_k)
             \left(-\tfrac12\log X\right)^{r+s-i-j}.
 \label{ct:logjets}
\end{align}
After specialization, only $k=m$ has zero exponent in $X$, and only
$i=r,j=s$ has logarithmic degree zero. Choose $K$ large enough that the
remainder and all fixed derivatives tend to zero. Equations
\eqref{ct:finite}, \eqref{ct:Cexp}, and \eqref{ct:Dexp} prove both formulas.
The degree bound in \eqref{ct:ck}--\eqref{ct:dk} proves the last assertion
about spectral derivatives. Fixed-cutoff argument derivatives follow from
the same uniform expansions.
\end{proof}

The arithmetic mean of the two alternating parity finite parts is the
endpoint continuation, including all its mixed jets. A full-sequence
ordinary cutoff constant is not supplied when the two constants differ.
Re-expanding in the parity index $(N-\epsilon)/2$ would be a different
coordinate convention. Similarly, imposing $q=-2m-p$ before differentiating
would calculate only tangent derivatives; it would lose the independent
transverse derivatives required by Theorem~\ref{ct:contacts}.

\subsection{Balanced coordinates, low-order contacts, and exact parity}

The most symmetric cutoff is
\begin{equation}\label{ct:balanced}
 X=N+\mu,\qquad \mu=\frac{a+b}{2}+1,\qquad
 \delta=\frac{a-b}{2}.
\end{equation}
The four centered arguments in the two endpoint products are then
$X+\delta-1/2$, $X-\delta+1/2$, $X+\delta+1/2$, and
$X-\delta-1/2$. Abbreviate the corresponding polynomials to $c_k,d_k$,
and put $S=p+q$, $D=p-q$. The first ambient coefficients are
\begin{align}
 c_0&=\tfrac12,&d_0&=0,\notag\\
 c_1&=-\frac{\delta D}{4},&d_1&=\frac D8,\label{ct:lowambient}\\
 c_2&=\frac{4\delta^2(D^2+2S)-D^2-2S^2}{64},&
 d_2&=-\frac{\delta(D^2+2S)}{16}.\notag
\end{align}
The third coefficients, useful for independent jets at the next resonance,
are
\begin{align}
 c_3&=-\frac{\delta D}{384}
 \{4\delta^2(D^2+6S+8)-3D^2-6S^2-36S\},\label{ct:thirdambient}\\
 d_3&=\frac D{768}
 \{12\delta^2(D^2+6S+8)-5D^2-6S^2-48S-16\}.\notag
\end{align}
For zero-order contacts only, their restrictions to $S=-2m$ give the
following compact table. Here $d=a-b=2\delta$ and $D=p-q$.
\begin{center}
\renewcommand{\arraystretch}{1.6}
\begin{tabular}{c|c|c}
$m$&symmetric contact $c_m$&alternating amplitude $d_m$\\\hline
$0$&$1/2$&$0$\\
$1$&$-dD/8$&$D/8$\\
$2$&$\displaystyle\frac{(D^2-8)d^2-D^2-32}{64}$
   &$\displaystyle-\frac{d(D^2-8)}{32}$\\
$3$&$\displaystyle-\frac{dD\{(D^2-28)d^2-3D^2\}}{768}$
   &$\displaystyle\frac{D\{3d^2(D^2-28)-5D^2+56\}}{768}$
\end{tabular}
\end{center}
The table must not replace the ambient polynomials when calculating
independent spectral derivatives.

\begin{proposition}[Balanced parity]\label{ct:parity}
At the balanced cutoff,
\begin{align}
 c_k(p,q;-\delta)&=(-1)^kc_k(p,q;\delta),&
 d_k(p,q;-\delta)&=(-1)^{k+1}d_k(p,q;\delta),\label{ct:reflection}\\
 c_k(q,p;\delta)&=(-1)^kc_k(p,q;\delta),&
 d_k(q,p;\delta)&=(-1)^kd_k(p,q;\delta).\notag
\end{align}
Every odd symmetric contact is divisible by $(a-b)(p-q)$.
Every even alternating contact vanishes at equal shifts, and every odd
alternating contact vanishes at equal orders.
\end{proposition}
\begin{proof}
Use $b_k(s,-v)=(-1)^kb_k(s,v)$ in the two endpoint products. Reflecting
$\delta$ exchanges these products; interchanging the independent orders
then gives the second line. Polynomial divisibility follows from the
resulting sign changes.
\end{proof}

The first off-diagonal symmetric contact is therefore
$(q-p)(a-b)/8$ on $p+q=-2$, and it vanishes on either diagonal separately.
For example,
\[
 p=\tfrac13,\quad q=-\tfrac73,\quad a=\tfrac23,\quad b=\tfrac75
 \quad\Longrightarrow\quad c_1=\frac{11}{45},\quad d_1=\frac13.
\]
At $p=q=s$, $a=b$, one has the exact reduction
\[
 C_N=2(N+a+1)^s\mathcal B_s(N+a+1)\mathcal B_s(N+a+2),
 \qquad D_N=0.
\]
Consequently the present coefficients specialize to twice the source's
diagonal endpoint coefficients. This checks all its even-order contacts,
all vanishing odd contacts, and every diagonal spectral derivative.

\subsection{Translation, the sharp derivative threshold, and moving shifts}

If $Y=X-h$, coefficient transport is finite:
\begin{equation}\label{ct:translate}
 c_m^{[h]}(p,q)=\sum_{k=0}^m c_k(p,q)
      \binom{-\lambda-k}{m-k}h^{m-k},
 \qquad
 d_m^{[h]}(p,q)=\sum_{k=0}^m d_k(p,q)
      \binom{-\lambda-k}{m-k}h^{m-k}.
\end{equation}
These are ambient polynomial identities, obtained by expanding
$X^{-\lambda-k}=Y^{-\lambda-k}(1+h/Y)^{-\lambda-k}$.
They retain total spectral degree at most $m$. At zero derivative order
and on $\lambda=-m$, each binomial factor is one. For higher independent
jets it must first be differentiated as a polynomial. Translating a cutoff
therefore changes transverse contact derivatives even when the restricted
zero-order formula appears elementary.

There is also a closed expression for every highest mixed derivative.
For a polynomial $P(p,q)$ write $[P]_m$ for its homogeneous component of
total spectral degree $m$.

\begin{theorem}[Highest mixed contacts and sharpness]\label{ct:top}
For the balanced cutoff translated by $h$, the formal generating identities
are
\begin{align}
 \sum_{m\ge0}[c_m^{[h]}]_m t^m
 &=\frac12\exp\left(-\frac{(p-q)\delta+(p+q)h}{2}t\right)
 \frac{\cosh((p-q)t/4)}{\cosh(pt/2)\cosh(qt/2)},\label{ct:topplus}\\
 \sum_{m\ge0}[d_m^{[h]}]_m t^m
 &=\frac12\exp\left(-\frac{(p-q)\delta+(p+q)h}{2}t\right)
 \frac{\sinh((p-q)t/4)}{\cosh(pt/2)\cosh(qt/2)}.\label{ct:topminus}
\end{align}
The total spectral degree is generically exactly $m$ for $c_m$, and
exactly $m$ for $d_m$ when $m\ge1$. Hence the threshold $r+s>m$ in
Theorem~\ref{ct:contacts} is sharp.
\end{theorem}
\begin{proof}
In \eqref{ct:Tformal}, the highest spectral coefficient of $t_j(s)$ is
$E_{2j}s^{2j}/(2^{2j}(2j)!)$: every nonconstant contribution from the
Gamma exponential has strictly smaller degree. The highest-degree
series for a shifted $\tau_s$ is therefore
$\tfrac12 e^{-svt/2}\operatorname{sech}(st/2)$. Multiplying the
appropriate two factors in \eqref{ct:ck}--\eqref{ct:dk}, and taking
their sum or difference, proves \eqref{ct:topplus}--\eqref{ct:topminus}.
For $h=q=0$, the coefficient of $\delta^mp^m$ in the symmetric
expression is $(-1)^m/(2^{m+1}m!)$. For $m\ge1$, the coefficient
of $\delta^{m-1}p^m$ in the alternating expression is
$(-1)^{m-1}/(2^{m+2}(m-1)!)$. Both are nonzero.
\end{proof}

The balanced coordinate depends on $a,b$. To differentiate the original
finite sums in these shifts, use \eqref{ct:ck}--\eqref{ct:dk} with fixed
$c$, differentiate, and only then set $c=\mu$. Holding the balanced $X$
artificially fixed while differentiating the original $N$-cutoff would
omit its coordinate derivative. This restriction does not affect
independent order derivatives, because $\mu$ is independent of $p,q$.

\subsection{Ordinary polynomial and logarithmic subtractions}

The coefficient algebra gives convergent sums representing every contact
and every jet. It does not require assigning ordinary values to divergent
unmodified sums. For $\Re c>1$ and $K\ge0$ define
\begin{equation}\label{ct:PK}
 P_K^+(p,q;X)=\sum_{k=0}^Kc_kX^{-\lambda-k},\qquad
 P_K^-(p,q;X)=\sum_{k=0}^Kd_kX^{-\lambda-k},\qquad X_n=n+c.
\end{equation}

\begin{theorem}[Normally convergent continuation]\label{ct:normal}
For $\Re\lambda>-K-1$,
\begin{equation}\label{ct:normalplus}
 \sum_{n\ge0}\{g_n^+-P_K^+(X_n-1)+P_K^+(X_n)\}
             =E_+-P_K^+(c-1)
\end{equation}
converges absolutely and locally normally. For $\Re\lambda>-K$,
\begin{equation}\label{ct:normalminus}
 \sum_{n\ge0}(-1)^n\{g_n^--P_K^-(X_n-1)-P_K^-(X_n)\}
             =E_--P_K^-(c-1)
\end{equation}
also converges absolutely and locally normally. The latter is an ordinary
convergent series, normally convergent after pairing consecutive terms,
on the larger domain $\Re\lambda>-K-1$.
All fixed mixed parameter derivatives may be taken under the sum in the
stated absolute-normal domains. Formula~\eqref{ct:logjets} specifies
every polynomial and logarithmic counterterm explicitly.
\end{theorem}
\begin{proof}
Subtract the finite telescoping identities for $P_K^\pm$ from
\eqref{ct:finite}. The remaining endpoints are
$C_N-P_K^+(X_N)$ and $D_N-P_K^-(X_N)$, both
$O(X_N^{-\Re\lambda-K-1})$. They vanish on the larger asserted domain.
For the symmetric series, the summand is a unit difference of its endpoint
remainder, so its size is $O(n^{-\Re\lambda-K-2})$. For the alternating
series the sum of two adjacent remainders is
$O(n^{-\Re\lambda-K-1})$, giving its absolute domain. Pairing successive
alternating terms turns this sum into a difference over two steps and
restores the extra inverse power. The uniform differentiated expansions
and Cauchy estimates justify all fixed derivatives in the stated domains.
The condition $\Re c>1$ keeps the lower power $c-1$ on the fixed branch.
\end{proof}

At $\lambda=-m$, the symmetric subtraction with $K=m$ is a polynomial.
Writing
\begin{equation}\label{ct:resonantP}
 \Pi_m^+(X)=\sum_{k=0}^mc_kX^{m-k},\qquad
 \Pi_m^-(X)=\sum_{k=0}^md_kX^{m-k},
\end{equation}
we obtain the ordinary absolutely convergent identity
\begin{equation}\label{ct:polynomialplus}
 \sum_{n\ge0}\{g_n^+-\Pi_m^+(n+\mu-1)+\Pi_m^+(n+\mu)\}
              =E_+-\Pi_m^+(\mu-1).
\end{equation}
Its summand is $O(n^{-2})$. The polynomial alternating subtraction is
ordinary convergent, generally conditionally, with value
$E_--\Pi_m^-(\mu-1)$. Including the additional term
$d_{m+1}X^{-1}$ makes it absolutely convergent. Independent mixed jets
at the resonance use $K=m$ in the symmetric formula and $K=m+1$ in
the alternating formula, with the logarithmic counterterms
\eqref{ct:logjets}. Thus the complete contact theorem has ordinary
convergent-series representatives at every derivative order.

\subsection{The first off-diagonal resonance and exact Gamma examples}

\begin{corollary}[A short identity on $p+q=-2$]\label{ct:firstresonance}
For $p+q=-2$ and $\Re a,\Re b>0$,
\begin{equation}\label{ct:firstplus}
 \boxed{\displaystyle
 \sum_{n\ge0}\left\{
  [V_p(n+a+1)+V_q(n+b+1)]u_nv_n+\frac12\right\}
 =\tau_p(a)\tau_q(b+1)+\tau_p(a+1)\tau_q(b)+\frac{pa+qb}{4}.}
\end{equation}
The series is absolutely and locally normally convergent on this
hyperplane. There is also the absolutely convergent alternating identity
\begin{align}
 &\sum_{n\ge0}(-1)^n\left\{
 [V_p(n+a+1)-V_q(n+b+1)]u_nv_n-\frac{p-q}{4}\right.\notag\\
 &\hspace{18mm}\left.
 -\frac{pq(a-b)}8\left[\frac1{n+(a+b)/2}
                      +\frac1{n+(a+b)/2+1}\right]\right\}\notag\\
 &\qquad=\tau_p(a)\tau_q(b+1)-\tau_p(a+1)\tau_q(b)
            -\frac{p-q}{8}-\frac{pq(a-b)}{4(a+b)}.
                                                    \label{ct:firstminus}
\end{align}
\end{corollary}
\begin{proof}
Here $\Pi_1^+(X)=X/2+(q-p)(a-b)/8$ and
\[
 \Pi_1^+(\mu-1)=-\frac{pa+qb}{4},
\]
where the relation $p+q=-2$ is used only after coefficient extraction.
Equation~\eqref{ct:polynomialplus} proves \eqref{ct:firstplus}.
For the alternating formula, \eqref{ct:lowambient} gives
$d_1=(p-q)/8$ and $d_2=pq(a-b)/8$ at this resonance. Apply
\eqref{ct:normalminus} with $P_2^-(X)=d_1+d_2/X$.
\end{proof}

In particular take $p=1$, $q=-3$, $a=1$, $b=2$, and set
\[
 u_n=\tau_1(n+2),\qquad v_n=\tau_{-3}(n+3).
\]
The two exact evaluations are
\begin{align}
 \sum_{n\ge0}\{[V_1(n+2)+V_{-3}(n+3)]u_nv_n+\tfrac12\}
 &=\frac{55\pi\log2}{128}
      +\frac{9(1-\log2)}{4\pi}-\frac54,\label{ct:exampleplus}\\
 \sum_{n\ge0}(-1)^n\left\{
 [V_1(n+2)-V_{-3}(n+3)]u_nv_n-1\right.&\notag\\[-3pt]
 \left.\hspace{6mm}-\frac38\left(\frac1{n+3/2}+\frac1{n+5/2}\right)
 \right\}
 &=\frac{55\pi\log2}{128}
      -\frac{9(1-\log2)}{4\pi}-\frac34.\label{ct:exampleminus}
\end{align}
Both sums are ordinary and absolutely convergent. Their summands use only
Gamma quotients, digamma functions, rational functions, and an Euler
polynomial: $\mathcal B_1(z)=\tfrac12\{\psi((z+1)/2)-\psi(z/2)\}$ and
\[
 \mathcal B_{-3}(z)=\frac12E_3(z),\qquad
 E_3(z)=z^3-\frac32z^2+\frac14.
\]
For a direct evaluation of the endpoints, put $r=\sqrt{2/\pi}$. Then
$h(1)=r$, $h(2)=1/r$, $h(3)=2r$,
$\mathcal B_1(1)=\log2$, and $\mathcal B_1(2)=1-\log2$. Consequently
\[
 \tau_1(1)\tau_{-3}(3)=\frac{55\pi\log2}{128},\qquad
 \tau_1(2)\tau_{-3}(2)=\frac{9(1-\log2)}{4\pi},
\]
which proves both displayed constants.

\subsection{Which resonant unmodified sums actually converge?}

The contact polynomials also resolve this ordinary-convergence question,
which is stronger than asking whether one constant coefficient vanishes.

\begin{proposition}[Ordinary convergence at negative-total-order resonances]
\label{ct:resonantconvergence}
On $p+q=-2m$ with $m>0$, the unmodified symmetric series does not converge.
At $m=0$ it converges absolutely and has the inherited contact $1/2$.
The alternating series at $m=0$ always converges; its convergence is
conditional unless $p=q=0$, when every summand is zero.
For the unmodified alternating series at $m\ge1$ the following are exact:
\begin{enumerate}
 \item If $p\ne q$, it does not converge.
 \item If $p=q=-m$ and $a\ne b$, it converges when $m=1$ and does
       not converge when $m\ge2$.
 \item If $p=q$ and $a=b$, every alternating summand is zero.
\end{enumerate}
In item 2 at $m=1$ the convergence is generally conditional.
\end{proposition}
\begin{proof}
The symmetric endpoint has leading term $X^m/2$, which cannot tend to a
finite limit for $m>0$. At $m=0$ it tends to $1/2$, and its difference
is $O(n^{-2})$. At $m=0$, the alternating endpoint tends to zero;
its summand has leading amplitude $(p-q)/(4n)$ unless $p=q=0$.
In that exceptional case $V_0=0$ and the summands vanish.
If $m\ge1$ and $p\ne q$, the alternating endpoint has leading term
$(p-q)X^{m-1}/8$, so $(-1)^ND_N$ has no limit. If $p=q=-m$, parity
and \eqref{ct:lowambient} instead give
\[
 D_N=\frac{m(a-b)}8X^{m-2}+O(X^{m-4}).
\]
For distinct shifts this tends to zero exactly when $m=1$; for $m=2$
it yields two different parity limits, and for larger $m$ it grows.
The exact finite identity \eqref{ct:finite} proves convergence or its
failure in each case. At $m=1$ the alternating summand generally has
nonzero $1/n$ amplitude, establishing the final qualification. Equal
orders and shifts give $u_n=v_n$, hence $D_n=0$ identically.
\end{proof}

The theorem answers the source's contact question at all its specified
hyperplanes. Independent extra outer powers and longer products remain
separate summation problems. The present finite coefficients offer a
precise starting point for those extensions: they identify which endpoint
terms must be retained before any additional spectral differentiation.

% END sections/02_contacts.tex


% BEGIN sections/03_stieltjes_dictionary.tex
\section{Stieltjes, zeta-derivative, and polygamma forms of the contacts}
\label{sj:section}

The independent-order theorem becomes an explicit identity in the
functions named in the research programme after a finite coefficient
conversion. The positive-order Stieltjes conversion itself is inherited
from \cite{ProveItBilinear}; its use with the completed negative-total-order
subtractions is the additional consequence here.

\subsection{A finite dictionary at every base order}

For a base order \(\sigma\) and an integer \(r\ge0\), write
\[
 \mathcal T_{\sigma,r}(z)
   =\left.\partial_s^r\tau_s(z)\right|_{s=\sigma},
 \qquad
 \mathcal V_{\sigma,r}(z)
   =\left.\partial_s^rV_s(z)\right|_{s=\sigma}.
\]
The formulas for \(\mathcal T\) use \(\Re z>0\); those for
\(\mathcal V\) use \(\Re z>1\).
Set \(L(z)=\ell(z)-\log2\), only in this section. At \(\sigma\ne1\),
\begin{equation}\label{sj:generalT}
 \mathcal T_{\sigma,r}(z)
 =\left(\frac{h(z)}2\right)^\sigma
  \sum_{j=0}^r\binom rj L(z)^{r-j}
  \left\{\zeta^{(j)}\!\left(\sigma,\frac z2\right)
        -\zeta^{(j)}\!\left(\sigma,\frac{z+1}{2}\right)\right\}.
\end{equation}
At the removable order \(\sigma=1\), the correct finite expression is
\begin{equation}\label{sj:stieltjesT}
 \mathcal T_{1,r}(z)=\frac{h(z)}2
  \sum_{j=0}^r\binom rj L(z)^{r-j}(-1)^j
  \left\{\gamma_j\!\left(\frac z2\right)
        -\gamma_j\!\left(\frac{z+1}{2}\right)\right\}.
\end{equation}
There is no pole or limiting quotient in this formula. At every \(\sigma\),
\begin{align}
 \mathcal V_{\sigma,r}(z)
 =h(z)^{-2\sigma}\bigl\{
 &z^\sigma(\log z-2\ell(z))^r \notag\\
 &-(z-1)^\sigma(\log(z-1)-2\ell(z))^r\bigr\}.
 \label{sj:Vjets}
\end{align}
Equations \eqref{sj:generalT}--\eqref{sj:Vjets} follow by differentiating
the exponential factor in \eqref{ct:B}--\eqref{ct:V}, and, at one, by
subtracting the two Stieltjes expansions before coefficient extraction.

For nonnegative integers \(k\),
\[
 \mathcal B_{-k}(z)=\frac12E_k(z),
 \qquad
 \frac{2e^{zt}}{e^t+1}=\sum_{k\ge0}E_k(z)\frac{t^k}{k!}.
\]
This replaces a zero-order negative-integer tail by an Euler polynomial.
An order derivative at \(-k\) is still given by
\eqref{sj:generalT}; it cannot be obtained by differentiating a polynomial
identity indexed only at integers.

\begin{corollary}[An explicit ordinary identity for every mixed boundary jet]
\label{sj:allboundary}
Fix \(\Re a,\Re b>0\), \(\sigma+\omega=-2m\), \(m\ge0\), and a cutoff \(c\) with
\(\Re c>1\). Let \(z_n=n+a+1\), \(w_n=n+b+1\), \(X_n=n+c\),
and use the fixed-cutoff coefficients \(c_k\) in \eqref{ct:ck}.
For \(r,s\ge0\) define
\begin{align}
 \mathcal G_{r,s}(z,w)
 ={}&\sum_{i=0}^r\binom ri
    \mathcal V_{\sigma,i}(z)
    \mathcal T_{\sigma,r-i}(z)\mathcal T_{\omega,s}(w)\notag\\
 &+\sum_{j=0}^s\binom sj
    \mathcal T_{\sigma,r}(z)\mathcal V_{\omega,j}(w)
    \mathcal T_{\omega,s-j}(w),\label{sj:G}\\
 \mathcal P_{r,s}(X)
 ={}&\sum_{k=0}^m X^{m-k}
   \sum_{i=0}^r\sum_{j=0}^s
       \binom ri\binom sj
       \left.\partial_p^i\partial_q^jc_k\right|_{(\sigma,\omega)}
       \left(-\frac{\log X}{2}\right)^{r+s-i-j}.
       \label{sj:P}
\end{align}
Then
\begin{align}
 &\sum_{n\ge0}
   \{\mathcal G_{r,s}(z_n,w_n)
          -\mathcal P_{r,s}(X_n-1)+\mathcal P_{r,s}(X_n)\}
       \notag\\
 &\quad=
 \mathcal T_{\sigma,r}(a)\mathcal T_{\omega,s}(b+1)
 +\mathcal T_{\sigma,r}(a+1)\mathcal T_{\omega,s}(b)
 -\mathcal P_{r,s}(c-1).
 \label{sj:ordinary}
\end{align}
The series is absolutely convergent. Its summand is
\(O(n^{-2}(1+\log^{r+s}n))\).
\end{corollary}

\begin{proof}
Take \(\partial_p^r\partial_q^s\) in
Theorem~\ref{ct:normal}, with \(K=m\), before specializing.
The point \((\sigma,\omega)\) lies inside its absolute-normal domain.
The product rule gives \eqref{sj:G}; differentiating the endpoint powers
gives \eqref{sj:P}. The differentiated remainder estimate proves the stated
decay and permits the interchange.
\end{proof}

Taking \(\sigma=1\), \(\omega=-2m-1\) in this corollary gives every
independent product of a Stieltjes order at half-shifted arguments with
negative-order Hurwitz derivatives. The correction to the cutoff
\emph{constant} vanishes whenever \(r+s>m\), by
Theorem~\ref{ct:contacts}. Polynomial and logarithmic subtraction terms may
still remain; their existence is not contradicted by the vanishing of that
single constant coefficient.

\subsection{An ordinary first-Stieltjes identity at total order zero}

The first independent derivative already gives a concrete formula whose
summands involve \(\gamma_1\), while its value uses only familiar constants.
For real \(z>1\), put
\begin{align*}
 T(z)&=\mathcal T_{1,1}(z),&
 R(z)&=\frac{z-\frac12}{2h(z)},\\
 A(z)&=h(z)^{-2}-\frac{h(z)^2}{z(z-1)},&
 W(z)&=\frac{z\log z-(z-1)\log(z-1)-2\ell(z)}{h(z)^2}.
\end{align*}
Here \(T\) is explicitly the two-term Stieltjes expression
\eqref{sj:stieltjesT}. Let
\[
 L=\log2,\qquad \lambda_1=\frac12\log\frac2\pi .
\]
Then
\begin{align}
 &\sum_{z=2}^{\infty}
   \left\{
      R(z)\bigl[W(z)\tau_1(z)+A(z)T(z)\bigr]
          -\frac14\log\frac z{z-1}
   \right\}\notag\\
 &\quad=
 \frac{3}{2\pi}\left\{L(\gamma+\lambda_1)-\frac{L^2}{2}\right\}
 +\frac{\pi}{8}\left\{-\lambda_1(1-L)-\gamma L+\frac{L^2}{2}\right\}.
 \label{sj:explicit}
\end{align}
This is an ordinary absolutely convergent series.

To prove it, use \eqref{sj:ordinary} at
\((\sigma,\omega,a,b,r,s)=(1,-1,1,1,1,0)\), with \(c=2\).
Its only endpoint subtraction is
\(\mathcal P_{1,0}(X)=-\frac14\log X\), giving precisely the logarithm
in \eqref{sj:explicit}. Now
\(\mathcal B_s(1)=\eta(s)=(1-2^{1-s})\zeta(s)\), and the removable
expansion at one gives
\[
 \eta(1)=L,\qquad \eta'(1)=\gamma L-\frac{L^2}{2}.
\]
Together with \(h(1)^2=2/\pi\), \(h(2)=1/h(1)\),
\(\mathcal B_s(2)=1-\eta(s)\), and
\(\mathcal B_{-1}(z)=(z-\frac12)/2\), this evaluates
\(T(1)R(2)+T(2)R(1)\) as the right side of \eqref{sj:explicit}.

\subsection{Argument derivatives and generalized harmonic values}

All argument derivatives of the dictionary reduce by the two laws
\begin{align}
 \partial_z^k\mathcal B_s(z)
    &=(-1)^k(s)_k\mathcal B_{s+k}(z),\label{sj:Barg}\\
 \ell^{(k)}(z)
    &=2^{-k}\left\{
       \psi^{(k-1)}\!\left(\frac{z+1}{2}\right)
       -\psi^{(k-1)}\!\left(\frac z2\right)\right\},
       \qquad k\ge1. \label{sj:ellarg}
\end{align}
The first follows by differentiating the Hurwitz series in its absolute
domain and continuing; the second follows directly from the Gamma quotient.
The product rule, or complete Bell polynomials for derivatives of an
exponential, therefore gives a terminating expression at every fixed
argument derivative order.

At integer arguments the classical recurrence gives
\[
 \psi(n)=H_{n-1}-\gamma,\qquad
 \psi^{(j)}(n)=(-1)^{j+1}j!
       \{\zeta(j+1)-H_{n-1}^{(j+1)}\}\quad(j\ge1).
\]
At half-integers, for \(n\ge0\) and \(j\ge1\) in the polygamma line,
\begin{align*}
 \psi(n+\tfrac12)&=-\gamma-2\log2+2H_{2n}-H_n,\\
 \psi^{(j)}(n+\tfrac12)
 &=(-1)^{j+1}j!\left\{
    (2^{j+1}-1)\zeta(j+1)-2^{j+1}H_{2n}^{(j+1)}
                               +H_n^{(j+1)}\right\}.
\end{align*}
Thus fixed integer and half-integer shift derivatives have an explicit
generalized harmonic alphabet. These elementary reductions are
classical \cite{DLMF}; the new boundary identities remain governed by the
ambient contact theorem, with the cutoff held fixed during shift
differentiation as specified there.

% END sections/03_stieltjes_dictionary.tex


% BEGIN sections/04_spectral.tex
\section{Scalar spectral jets and ordinary resolvent integrals}
\label{sp:section}

Throughout this section $a>0$, $x_n=n+a$ for $n=0,1,\ldots$, and
$\mu_1,\ldots,\mu_N$ are positive real numbers, with $N\geq1$.
The logarithm $\Log$ is the principal logarithm on
$\mathbb C\setminus(-\infty,0]$.  Powers of positive real numbers use
their real logarithms.  We write $\zeta(s,a)$ for the Hurwitz zeta
function, $\psi$ for the digamma function, and
$\zeta^{(r)}(s,a)=\partial_s^r\zeta(s,a)$.
The formulas below are specialized deductions for scalar factors and
their connected logarithmic products.  They make no assertion of
global priority for the general principle relating spectral finite
parts and local asymptotic data.

\subsection{The scalar master formula}
\label{sp:scalar-subsection}

For $\Re s>1/(2N)$ define
\begin{equation}
 Z_N(s;a,\boldsymbol\mu)
 :=\sum_{n=0}^{\infty}\prod_{j=1}^{N}(x_n^2+\mu_j)^{-s}.
 \label{sp:zeta-definition}
\end{equation}
For integers $r\geq1$, its $r$th spectral jet is
$J_{N,r}(a;\boldsymbol\mu):=\partial_s^rZ_N(0;a,\boldsymbol\mu)$,
once continuation to a neighborhood of zero has been established.

\begin{lemma}[Continuation by a normally convergent difference]
\label{sp:continuation}
The function in \eqref{sp:zeta-definition} has the continuation
\begin{equation}
 Z_N(s;a,\boldsymbol\mu)
 =\zeta(2Ns,a)+
 \sum_{n=0}^{\infty}x_n^{-2Ns}
 \left\{\exp\left[-s\sum_{j=1}^{N}
             \log\left(1+\frac{\mu_j}{x_n^2}\right)\right]-1\right\},
 \label{sp:normal-continuation}
\end{equation}
where the series is normally convergent on $\Re s>-1/(2N)$.
In particular $Z_N$ is holomorphic near $s=0$, $Z_N(0)=\zeta(0,a)$,
and
\begin{equation}
 J_{N,r}=(2N)^r\zeta^{(r)}(0,a)
          +(-1)^r\sum_{n=0}^{\infty}F_{N,r}(x_n^2),
 \qquad
 F_{N,r}(z):=
 \left(\sum_{j=1}^{N}\Log(z+\mu_j)\right)^r-N^r\Log^r z.
 \label{sp:jet-ordinary-series}
\end{equation}
The series in \eqref{sp:jet-ordinary-series} is absolutely convergent.
\end{lemma}

\begin{proof}
Put $L_n=\sum_j\log(1+\mu_j/x_n^2)$.  Then $L_n=O(x_n^{-2})$.
For a compact subset $K$ of $\Re s>-1/(2N)$, the tail of the
series in \eqref{sp:normal-continuation} is bounded, uniformly on $K$,
by a constant times $x_n^{-2N\inf_K\Re s-2}$.  This is summable.
The finitely many remaining summands are entire functions of $s$.
The identity first holds in the defining half-plane by factoring
out $x_n^{-2Ns}$, and hence gives the asserted continuation.
The sole pole of $\zeta(2Ns,a)$ in this region is away from zero.
Normal convergence permits differentiation at zero, giving
\eqref{sp:jet-ordinary-series}.  Finally,
\[
 \sum_j\log(x_n^2+\mu_j)=N\log(x_n^2)+O(x_n^{-2}),
 \qquad
 F_{N,r}(x_n^2)=O\left(\frac{1+\log^{r-1}x_n}{x_n^2}\right)
\]
for large $n$, which also proves absolute convergence directly.
\end{proof}

For $t>0$ distinct from the masses, set
\begin{equation}
 A(t):=\sum_{j=1}^{N}\log|\mu_j-t|,
 \qquad k(t):=\#\{j:\mu_j<t\},
 \label{sp:cut-data}
\end{equation}
counting repeated masses with their multiplicity, and define the
real, possibly signed density
\begin{equation}
 \rho_{N,r}(t):=-\frac{1}{\pi}\Im\left\{
        [A(t)+i\pi k(t)]^r
           -N^r[\log t+i\pi]^r\right\}.
 \label{sp:scalar-density}
\end{equation}
Its values at the finitely many mass points may be assigned
arbitrarily.  They do not change any integral below.

\begin{theorem}[A single ordinary integral for every scalar jet]
\label{sp:master}
For every integer $N\geq1$, every integer $r\geq1$, every $a>0$,
and all positive masses,
\begin{equation}
 \boxed{\displaystyle
 J_{N,r}(a;\boldsymbol\mu)
 =(2N)^r\zeta^{(r)}(0,a)
       +(-1)^r\int_0^{\infty}\rho_{N,r}(t)Q_a(t)\,dt,}
 \label{sp:master-formula}
\end{equation}
where
\begin{equation}
 Q_a(t):=\sum_{n=0}^{\infty}\frac{1}{x_n^2+t}
 =\frac{\psi(a+i\sqrt t)-\psi(a-i\sqrt t)}{2i\sqrt t},
 \qquad Q_a(0):=\psi'(a).
 \label{sp:resolvent-kernel}
\end{equation}
The integral is absolutely convergent.  At the mass points it is
understood as the sum of the ordinary improper integrals over the
intervals separated by the distinct masses; no finite part or
principal value is required.
\end{theorem}

\begin{proof}
We first establish the Stieltjes representation
\begin{equation}
 F_{N,r}(z)=\int_0^{\infty}\frac{\rho_{N,r}(t)}{z+t}\,dt,
 \qquad z\in\mathbb C\setminus(-\infty,0].
 \label{sp:scalar-stieltjes}
\end{equation}
The function $F_{N,r}$ is holomorphic on this slit plane and
satisfies Schwarz reflection.  On the upper bank of the cut,
\[
 \sum_j\Log(\mu_j-t+i0)=A(t)+i\pi k(t),
 \qquad \Log(-t+i0)=\log t+i\pi.
\]
Consequently its upper boundary value has imaginary part
$-\pi\rho_{N,r}(t)$.  Near zero,
$F_{N,r}(z)=O(1+|\Log z|^r)$; at a nonzero endpoint $-\mu_j$
it has at worst a power of a logarithm.  The integrals on circles
of radius $\varepsilon$ around these endpoints therefore tend to
zero, since $\varepsilon(1+|\log\varepsilon|^r)\to0$.
At infinity, uniformly along a large keyhole contour,
\[
 \sum_j\Log(z+\mu_j)
       =N\Log z+O(|z|^{-1}),
 \qquad
 F_{N,r}(z)=O\left(\frac{1+\log^{r-1}|z|}{|z|}\right).
\]
The outer circle in the Cauchy integral hence contributes zero
in the limit.  Cauchy's formula, followed by collapsing the two
banks of the cut, gives \eqref{sp:scalar-stieltjes} with the sign
in \eqref{sp:scalar-density}.

Here are explicit bounds that justify the resulting improper
integral and its later interchange with a series.  Since the
masses are positive, $k(t)=0$ for sufficiently small $t$ and
$A(t)$ remains bounded there.  Thus
\begin{equation}
 \rho_{N,r}(t)=O(1+|\log t|^{r-1})\qquad(t\downarrow0).
 \label{sp:density-zero-bound}
\end{equation}
At each mass point the density is bounded by a constant times
$1+|\log|t-\mu_j||^{r-1}$, allowing for multiplicities; these
singularities are locally integrable.  Beyond the largest mass,
write
\[
 A(t)=N\log t+\delta(t),\qquad
 \delta(t)=\sum_j\log(1-\mu_j/t)=O(t^{-1}),\qquad k(t)=N.
\]
In the difference
$[N(\log t+i\pi)+\delta(t)]^r-[N(\log t+i\pi)]^r$,
taking the imaginary part lowers the maximal logarithmic degree
by one.  Hence, for $r\geq2$,
\begin{equation}
 \rho_{N,r}(t)=O\left(\frac{1+\log^{r-2}t}{t}\right)
       \qquad(t\to\infty).
 \label{sp:density-infinity-bound}
\end{equation}
For $r=1$ the density is zero beyond the largest mass.
These estimates also directly give absolute convergence in
\eqref{sp:scalar-stieltjes} at any fixed $z$ off the cut.

The convergent digamma difference series gives
\eqref{sp:resolvent-kernel}.  It also shows that $Q_a(t)>0$,
that $Q_a$ is continuous at zero with value $\psi'(a)$, and that
it is bounded on compact intervals of $t\geq0$.
Comparison of its decreasing summand with an integral yields
\[
 Q_a(t)=\int_a^{\infty}\frac{dx}{x^2+t}+O(t^{-1})
       =\frac{\pi}{2\sqrt t}+O(t^{-1})
       \qquad(t\to\infty).
\]
It follows from the endpoint bounds that
\[
 \int_0^{\infty}|\rho_{N,r}(t)|Q_a(t)\,dt<\infty.
\]
Tonelli's theorem applied to the absolute values, followed by
Fubini's theorem, now proves
\[
 \sum_{n=0}^{\infty}F_{N,r}(x_n^2)
   =\int_0^{\infty}\rho_{N,r}(t)
                     \sum_{n=0}^{\infty}\frac{dt}{x_n^2+t}
   =\int_0^{\infty}\rho_{N,r}(t)Q_a(t)\,dt.
\]
Substitution in \eqref{sp:jet-ordinary-series} gives the result.
\end{proof}

\begin{remark}[Explicit densities and collisions]
\label{sp:explicit-densities}
Expanding only the odd binomial terms gives, at every order,
\begin{equation}
 \rho_{N,r}(t)=
 \sum_{h=0}^{\lfloor(r-1)/2\rfloor}
   (-1)^h\binom{r}{2h+1}\pi^{2h}
   \left\{N^r(\log t)^{r-2h-1}
                -k(t)^{2h+1}A(t)^{r-2h-1}\right\}.
 \label{sp:binomial-density}
\end{equation}
In particular,
\begin{align}
 \rho_{N,1}(t)&=N-k(t),\label{sp:density-one}\\
 \rho_{N,2}(t)&=2N^2\log t-2k(t)A(t),\label{sp:density-two}\\
 \rho_{N,3}(t)&=3N^3\log^2t-3k(t)A(t)^2
                    +\pi^2\{k(t)^3-N^3\}.
 \label{sp:density-three}
\end{align}
The densities need not be nonnegative at higher orders.
Repeated masses are included directly in these formulas: $k$
jumps by their multiplicity, and $A$ contains the corresponding
multiple of $\log|\mu-t|$.  The proof of
Theorem~\ref{sp:master} applies without a distinctness assumption,
so no limiting prescription is needed when masses coincide.
\end{remark}

\begin{corollary}[The first jet and its normalization]
\label{sp:first-jet}
Let $\log\Gamma$ denote the analytic branch on $\Re z>0$ that
is real on the positive real axis.  Then
\begin{equation}
 J_{N,1}(a;\boldsymbol\mu)
 =\sum_{j=1}^{N}\left\{
       \log\Gamma(a+i\sqrt{\mu_j})
       +\log\Gamma(a-i\sqrt{\mu_j})-\log(2\pi)\right\}.
 \label{sp:first-jet-gamma}
\end{equation}
\end{corollary}

\begin{proof}
Since $N-k(t)=\sum_j\boldsymbol1_{(0,\mu_j)}(t)$, the integral
in \eqref{sp:master-formula} is a sum of $\int_0^{\mu_j}Q_a(t)\,dt$.
Differentiation with respect to $t$ gives
\[
 \frac{d}{dt}\left\{
   \log\Gamma(a+i\sqrt t)+\log\Gamma(a-i\sqrt t)\right\}
 =-Q_a(t).
\]
The identity extends continuously to $t=0$.  Therefore
\[
 \int_0^\mu Q_a(t)\,dt
 =2\log\Gamma(a)-\log\Gamma(a+i\sqrt\mu)
                       -\log\Gamma(a-i\sqrt\mu).
\]
Combining this with
$\zeta'(0,a)=\log\Gamma(a)-\tfrac12\log(2\pi)$ proves
\eqref{sp:first-jet-gamma}.
\end{proof}

\subsection{Parameter calculus without regularized interchange}
\label{sp:parameters}

\begin{proposition}[Derivatives and an anchored primitive in $a$]
\label{sp:a-calculus}
For every nonnegative integer $m$,
\begin{equation}
 \partial_a^mJ_{N,r}
 =(2N)^r\partial_a^m\zeta^{(r)}(0,a)
   +(-1)^r\int_0^{\infty}\rho_{N,r}(t)Q_a^{[m]}(t)\,dt,
 \label{sp:a-derivatives}
\end{equation}
where, with $\psi^{(0)}=\psi$,
\begin{equation}
 Q_a^{[m]}(t):=
 \frac{\psi^{(m)}(a+i\sqrt t)-\psi^{(m)}(a-i\sqrt t)}
      {2i\sqrt t},
 \qquad Q_a^{[m]}(0)=\psi^{(m+1)}(a).
 \label{sp:polygamma-kernel}
\end{equation}
All these integrals are absolutely convergent and locally uniform
for $a>0$.

For a fixed anchor $a_*>0$, define
\begin{align}
 K_a(t)&:=\frac{\log\Gamma(a+i\sqrt t)-\log\Gamma(a-i\sqrt t)}
                      {2i\sqrt t},\qquad t>0,\notag\\
 H_{a,a_*}(t)&:=K_a(t)-K_{a_*}(t),
 \label{sp:gamma-kernel}\\
 H_{a,a_*}(0)&:=\psi(a)-\psi(a_*),\notag\\
 B_{N,r}(a)
 &:=(2N)^r r!\sum_{j=0}^{r}\frac{\zeta^{(j)}(-1,a)}{j!}.
 \label{sp:primitive-baseline}
\end{align}
The branch of $\log\Gamma$ is the one used in
Corollary~\ref{sp:first-jet}.  Then the anchored primitive is
\begin{equation}
 \int_{a_*}^{a}J_{N,r}(u;\boldsymbol\mu)\,du
 =B_{N,r}(a)-B_{N,r}(a_*)
     +(-1)^r\int_0^{\infty}
                      \rho_{N,r}(t)H_{a,a_*}(t)\,dt.
 \label{sp:anchored-primitive}
\end{equation}
The final integral is again ordinary and absolutely convergent.
\end{proposition}

\begin{proof}
Let $a$ range over a compact interval $I\subset(0,\infty)$.
Choose a positive complex radius $\eta$ smaller than one half of
$\inf I$.  For real $a\in I$, $t\geq0$, and $|w-a|\leq\eta$,
each factor $n+w\pm i\sqrt t$ has modulus at least a fixed
positive multiple of $\sqrt{(n+a)^2+t}$.  Cauchy's derivative
estimate applied to $((n+w)^2+t)^{-1}$ therefore gives
\[
 \left|\partial_a^m\frac{1}{(n+a)^2+t}\right|
 \leq\frac{C_{I,m}}{(n+a)^2+t}.
\]
Summing and using the bounds in the proof of
Theorem~\ref{sp:master} produces a majorant for every derivative
that is integrable against $|\rho_{N,r}(t)|\,dt$, uniformly for
$a\in I$.  Differentiation under the integral is valid to every
order.  Differentiating the digamma expression for $Q_a$ gives
\eqref{sp:polygamma-kernel}; its value at zero follows by a Taylor
expansion.  The same local uniform estimates justify derivatives
of the normally convergent series in
\eqref{sp:normal-continuation}.

For the primitive, integration of $Q_u(t)$ in $u$ gives
\[
 H_{a,a_*}(t)=\int_{a_*}^{a}Q_u(t)\,du.
\]
This identity both fixes the branch normalization and proves the
required integral bound, including the continuous value at $t=0$.
For the Hurwitz term, the identity
$\partial_a\zeta(v,a)=-v\zeta(v+1,a)$ gives
\[
 \partial_a\left\{-\frac{\zeta(2Ns-1,a)}{2Ns-1}\right\}
 =\zeta(2Ns,a).
\]
Writing $u=2Ns$ and expanding at $u=0$,
\[
 -\frac{\zeta(2Ns-1,a)}{2Ns-1}
 =\frac{\zeta(u-1,a)}{1-u},\qquad
 \left.\partial_s^r
    \frac{\zeta(2Ns-1,a)}{1-2Ns}\right|_{s=0}
 =(2N)^r r!\sum_{j=0}^{r}\frac{\zeta^{(j)}(-1,a)}{j!}.
\]
Thus $B_{N,r}'(a)=(2N)^r\zeta^{(r)}(0,a)$.
Integrating \eqref{sp:master-formula}, with the locally uniform
absolute bounds just proved, gives
\eqref{sp:anchored-primitive}.
\end{proof}

\begin{proposition}[Mixed mass derivatives and resolvents]
\label{sp:mass-derivatives}
Hold $N$ fixed.  Choose $q$ distinct labels
$i_1,\ldots,i_q\in\{1,\ldots,N\}$ and put
$S_n:=\sum_{j=1}^{N}\log(x_n^2+\mu_j)$.
For $1\leq q\leq r$,
\begin{equation}
 \partial_{\mu_{i_1}}\cdots\partial_{\mu_{i_q}}J_{N,r}
 =(-1)^r\frac{r!}{(r-q)!}
   \sum_{n=0}^{\infty}
     \frac{S_n^{r-q}}{\prod_{\ell=1}^{q}(x_n^2+\mu_{i_\ell})}.
 \label{sp:mixed-mass-formula}
\end{equation}
For $q>r$ the mixed derivative vanishes.  Every series in
\eqref{sp:mixed-mass-formula} is absolutely convergent, including
when different labels carry equal masses.
In particular, if $r\leq N$, then
\begin{equation}
 \partial_{\mu_{i_1}}\cdots\partial_{\mu_{i_r}}J_{N,r}
 =(-1)^r r!\sum_{n=0}^{\infty}
          \frac{1}{\prod_{\ell=1}^{r}(x_n^2+\mu_{i_\ell})}
 =-r!\,Q_a[\mu_{i_1},\ldots,\mu_{i_r}],
 \label{sp:rational-mass-resolvent}
\end{equation}
where brackets denote the ordinary divided difference of order
$r-1$ with its confluent continuation.
\end{proposition}

\begin{proof}
In a common domain of absolute convergence, differentiating the
defining spectral summand in $q$ distinct mass labels gives
\[
 \partial_{\mu_{i_1}}\cdots\partial_{\mu_{i_q}}Z_N(s)
 =(-s)^q\sum_{n=0}^{\infty}
     \frac{e^{-sS_n}}{\prod_{\ell=1}^{q}(x_n^2+\mu_{i_\ell})}.
\]
For $q\geq1$ the series on the right is normally convergent in
a neighborhood of $s=0$, uniformly for masses in compact subsets
of the positive orthant.  It supplies the continuation of the
mass derivative.  Taking its $r$th derivative at zero proves
\eqref{sp:mixed-mass-formula} and the vanishing for $q>r$.
The summands are bounded by a constant times
$x_n^{-2q}(1+\log^{r-q}x_n)$ when $q\leq r$.
Finally,
\[
 \left[t\longmapsto\frac{1}{x^2+t}\right]
        [\nu_1,\ldots,\nu_r]
 =\frac{(-1)^{r-1}}{\prod_{j=1}^{r}(x^2+\nu_j)}.
\]
Summing this identity proves the divided-difference formula.
Normal convergence in the masses justifies confluence at equal
arguments.
\end{proof}

\begin{remark}[Repeated differentiation of one mass]
\label{sp:repeated-mass-derivatives}
The qualification ``distinct labels'' in
Proposition~\ref{sp:mass-derivatives} concerns the differentiation
variables, not their numerical values.  A formula for an arbitrary
nonzero multi-index $\boldsymbol\alpha=(\alpha_1,\ldots,\alpha_N)$
is obtained without differentiating a moving branch-cut density.
Let $d=\sum_j\alpha_j$, let $q=\#\{j:\alpha_j>0\}$, and write
\[
 \prod_{j=1}^{N}(s)_{\alpha_j}
      =\sum_{\ell=q}^{d}c_{\boldsymbol\alpha,\ell}s^\ell,
 \qquad (s)_m=s(s+1)\cdots(s+m-1),\quad (s)_0=1.
\]
Then
\begin{equation}
 \partial_{\boldsymbol\mu}^{\boldsymbol\alpha}J_{N,r}
 =(-1)^d r!\sum_{\ell=q}^{\min(d,r)}
       \frac{(-1)^{r-\ell}c_{\boldsymbol\alpha,\ell}}{(r-\ell)!}
       \sum_{n=0}^{\infty}
       \frac{S_n^{r-\ell}}
            {\prod_{j=1}^{N}(x_n^2+\mu_j)^{\alpha_j}}.
 \label{sp:all-mass-derivatives}
\end{equation}
An empty outer sum is zero.  Indeed, differentiating the spectral
summand gives
\[
 (-1)^d\prod_j(s)_{\alpha_j}\,e^{-sS_n}
                  \prod_j(x_n^2+\mu_j)^{-\alpha_j},
\]
whose series is normal near
zero; coefficient extraction gives
\eqref{sp:all-mass-derivatives}.  Each inner series is ordinary
and absolutely convergent.
\end{remark}

\begin{corollary}[Integral shifts and the zero-mass boundary]
\label{sp:integer-a}
For an integer $M\geq1$ and $t>0$,
\begin{equation}
 Q_M(t)=\frac{\pi\sqrt t\,\coth(\pi\sqrt t)-1}{2t}
           -\sum_{j=1}^{M-1}\frac{1}{j^2+t}.
 \label{sp:integer-resolvent}
\end{equation}
For integers $p\geq1$ its confluent derivatives give
\begin{equation}
 \sum_{n=0}^{\infty}\frac{1}{((n+M)^2+t)^p}
 =\frac{(-1)^{p-1}}{(p-1)!}\partial_t^{p-1}Q_M(t).
 \label{sp:confluent-resolvent}
\end{equation}
These expressions extend continuously to $t=0$, where
\begin{equation}
 \sum_{n=0}^{\infty}(n+M)^{-2p}
    =\zeta(2p)-H_{M-1}^{(2p)},
 \qquad H_m^{(v)}:=\sum_{j=1}^{m}j^{-v},\quad H_0^{(v)}:=0.
 \label{sp:harmonic-resolvent}
\end{equation}
Consequently, when $r\leq N$, the limit of
\eqref{sp:rational-mass-resolvent} as its $r$ selected masses tend
to zero is
\begin{equation}
 (-1)^r r!\{\zeta(2r)-H_{M-1}^{(2r)}\},\qquad a=M.
 \label{sp:harmonic-mixed-jet}
\end{equation}
The remaining masses do not enter this mixed derivative.
\end{corollary}

\begin{proof}
The partial-fraction expansion of the hyperbolic cotangent gives
\[
 \sum_{j=1}^{\infty}(j^2+t)^{-1}
       =\frac{\pi\sqrt t\coth(\pi\sqrt t)-1}{2t}.
\]
Removing its first $M-1$ terms proves
\eqref{sp:integer-resolvent}.  Normal convergence permits the
$t$ derivatives in \eqref{sp:confluent-resolvent}.  At $t=0$ the
series becomes a tail of the defining series for $\zeta(2p)$,
which proves \eqref{sp:harmonic-resolvent}.  Dominated convergence
in \eqref{sp:rational-mass-resolvent} gives the last assertion.
\end{proof}

\subsection{Compact densities for connected logarithmic products}
\label{sp:connected-subsection}

We now use $r$ for the number of logarithmic factors.  Fix an
anchor $\beta\geq0$, let $\nu_1,\ldots,\nu_r>\beta$, and put
$M=\max_j\nu_j$.  Define
\begin{equation}
 G_{\boldsymbol\nu,\beta}(z)
 :=\prod_{j=1}^{r}
       \log\frac{z+\nu_j}{z+\beta}.
 \label{sp:connected-product}
\end{equation}
Each logarithm means the analytic continuation from positive
real $z$ of $\Log(z+\nu_j)-\Log(z+\beta)$, with value tending
to zero at infinity.  This gives a holomorphic function on
$\mathbb C\setminus[-M,-\beta]$.
For $\beta<t<M$ away from the masses, write
\begin{equation}
 \ell_j(t):=\log\left|\frac{\nu_j-t}{\beta-t}\right|,
 \qquad \chi_j(t):=\boldsymbol1_{\{t<\nu_j\}},
 \qquad
 \rho_{\boldsymbol\nu,\beta}(t)
 :=-\frac1\pi\Im\prod_{j=1}^{r}
                     \{\ell_j(t)-i\pi\chi_j(t)\}.
 \label{sp:connected-density}
\end{equation}
The density is taken to be zero outside $[\beta,M]$.

\begin{theorem}[Compact representation and exact moments]
\label{sp:compact-density}
The density in \eqref{sp:connected-density} is integrable and
satisfies
\begin{equation}
 G_{\boldsymbol\nu,\beta}(z)
 =\int_\beta^M\frac{\rho_{\boldsymbol\nu,\beta}(t)}{z+t}\,dt,
 \qquad z\in\mathbb C\setminus[-M,-\beta].
 \label{sp:compact-stieltjes}
\end{equation}
For every integer $k\geq0$,
\begin{equation}
 \int_\beta^M t^k\rho_{\boldsymbol\nu,\beta}(t)\,dt
 =(-1)^k[u^{k+1}]
       \prod_{j=1}^{r}\log\frac{1+\nu_j u}{1+\beta u},
 \label{sp:compact-moments}
\end{equation}
where $[u^m]$ extracts a Taylor coefficient at $u=0$.
In particular, with $C:=\prod_{j=1}^{r}(\nu_j-\beta)$,
\begin{align}
 \int_\beta^M t^k\rho_{\boldsymbol\nu,\beta}(t)\,dt&=0,
       &&0\leq k\leq r-2,\label{sp:vanishing-moments}\\
 \int_\beta^M t^{r-1}\rho_{\boldsymbol\nu,\beta}(t)\,dt
      &=(-1)^{r-1}C.
 \label{sp:first-nonzero-moment}
\end{align}
The first range is empty if $r=1$.
\end{theorem}

\begin{proof}
On the upper cut $z=-t+i0$, the denominator logarithm has
imaginary part $\pi$ for $t>\beta$, while the numerator
logarithm has imaginary part $\pi$ precisely for $t>\nu_j$.
Their difference is therefore
$\ell_j(t)-i\pi\chi_j(t)$.
For $t>M$ all these differences have identical real boundary
values on both banks of the cut, which explains the compact
support.  At every endpoint the product has at worst a power of
a logarithm.  Thus the density is integrable, even when several
masses coincide, and every small endpoint circle contributes
$O(\varepsilon(1+|\log\varepsilon|^r))$ to the Cauchy integral.
At infinity,
\begin{equation}
 G_{\boldsymbol\nu,\beta}(z)=Cz^{-r}+O(z^{-r-1}).
 \label{sp:connected-asymptotic}
\end{equation}
The outer contour therefore vanishes.  Collapsing the Cauchy
contour onto the finite cut proves
\eqref{sp:compact-stieltjes}, with its indicated sign.

For $|z|>M$, expand $(z+t)^{-1}$ into its uniformly convergent
geometric series on the support of the density.  The coefficient
of $z^{-k-1}$ in the resulting integral is
$(-1)^k\int t^k\rho_{\boldsymbol\nu,\beta}(t)\,dt$.
On the other hand, the same coefficient on the left is
$[u^{k+1}]\prod_j\log((1+\nu_j u)/(1+\beta u))$.
This proves \eqref{sp:compact-moments}.  Each logarithmic factor
starts with $(\nu_j-\beta)u$, so the product starts with $Cu^r$.
The final two assertions follow.
\end{proof}

\begin{corollary}[Coincident masses]
\label{sp:coincident-density}
If $\nu_1=\cdots=\nu_r=\nu>\beta$, set
$y=\log((\nu-t)/(t-\beta))$ for $\beta<t<\nu$.  Then
\begin{equation}
 \rho_{\nu,\beta}^{(r)}(t)=P_r(y),\qquad
 P_r(y):=\sum_{h=0}^{\lfloor(r-1)/2\rfloor}
       (-1)^h\binom{r}{2h+1}\pi^{2h}y^{r-2h-1}.
 \label{sp:coincident-polynomial}
\end{equation}
For example,
\begin{equation}
 P_1(y)=1,\qquad P_2(y)=2y,\qquad
 P_3(y)=3y^2-\pi^2,\qquad P_4(y)=4y^3-4\pi^2y.
 \label{sp:coincident-examples}
\end{equation}
\end{corollary}

\begin{proof}
On this single interval every factor in
\eqref{sp:connected-density} equals $y-i\pi$.
Expanding $-\pi^{-1}\Im(y-i\pi)^r$ gives the formula.
\end{proof}

\begin{corollary}[An ordinary connected spectral sum]
\label{sp:ordinary-connected}
For $a>0$,
\begin{equation}
 \sum_{n=0}^{\infty}
   \prod_{j=1}^{r}\log\frac{x_n^2+\nu_j}{x_n^2+\beta}
 =\int_\beta^M\rho_{\boldsymbol\nu,\beta}(t)Q_a(t)\,dt.
 \label{sp:connected-sum}
\end{equation}
Both sides are ordinary and absolutely convergent.
\end{corollary}

\begin{proof}
The summand is $O(x_n^{-2r})$ by
\eqref{sp:connected-asymptotic}.  Moreover, $Q_a$ is bounded on
the compact interval $[\beta,M]$, including when $\beta=0$.
Thus $\int|\rho_{\boldsymbol\nu,\beta}(t)|Q_a(t)\,dt$ is finite,
and absolute Fubini applied to
\eqref{sp:compact-stieltjes} proves the assertion.
\end{proof}

\begin{remark}[A fixed number of spectral factors]
\label{sp:fixed-N-differences}
Suppose $N\geq r$ is fixed, and choose $r$ labeled mass slots.
In slot $j$, let $\Delta_j$ mean the value with mass $\nu_j$
minus the value with mass $\beta$; the other slots are held
fixed.  The normally convergent formula
\eqref{sp:normal-continuation} also permits zero masses at this
fixed positive $a$.  Consequently,
\begin{equation}
 \Delta_1\cdots\Delta_r J_{N,r}
 =(-1)^r r!\int_\beta^M
                    \rho_{\boldsymbol\nu,\beta}(t)Q_a(t)\,dt.
 \label{sp:connected-finite-difference}
\end{equation}
Indeed, in \eqref{sp:jet-ordinary-series} the mass-independent
baseline cancels, while the mixed difference of the $r$th power
of a sum of logarithms is $r!$ times the product of their $r$
increments.  This is a replacement operation at fixed $N$.
Adding or deleting a spectral factor changes the scale $2N$ in
the zeta term and is a different operation.
\end{remark}

\begin{theorem}[The critical odd weight as an ordinary sum]
\label{sp:critical-weight}
Under the assumptions of Theorem~\ref{sp:compact-density},
\begin{equation}
 \boxed{\begin{aligned}
 &\sum_{n=0}^{\infty}\left\{
        x_n^{2r-1}G_{\boldsymbol\nu,\beta}(x_n^2)-\frac{C}{x_n}
                         \right\}\\
 &\qquad=(-1)^r\int_\beta^M
       \rho_{\boldsymbol\nu,\beta}(t)t^{r-1}
       \{\Re\psi(a+i\sqrt t)-\psi(a)\}\,dt.
 \end{aligned}}
 \label{sp:critical-ordinary-sum}
\end{equation}
The summand on the left is $O(x_n^{-3})$.  The integral on the
right is an ordinary absolutely convergent integral; the
digamma difference has its continuous value zero at $t=0$.
\end{theorem}

\begin{proof}
Polynomial division gives
\[
 \frac{x^{2r-1}}{x^2+t}
 =\sum_{k=0}^{r-2}(-t)^k x^{2r-3-2k}
          +(-t)^{r-1}\frac{x}{x^2+t},
\]
where the sum is empty for $r=1$.
Upon integration against $\rho_{\boldsymbol\nu,\beta}$, every
polynomial term vanishes by \eqref{sp:vanishing-moments}.
Using \eqref{sp:first-nonzero-moment} for the subtraction gives
the exact identity, for each $x>0$,
\[
 x^{2r-1}G_{\boldsymbol\nu,\beta}(x^2)-\frac{C}{x}
 =(-1)^{r-1}\int_\beta^M
       \rho_{\boldsymbol\nu,\beta}(t)t^{r-1}
             \left\{\frac{x}{x^2+t}-\frac1x\right\}\,dt.
\]
The expression in braces equals $-t/[x(x^2+t)]$ and is bounded
in absolute value by $M/x^3$ on the whole integration interval.
Hence absolute Fubini is valid after this subtraction.
The convergent digamma difference series yields
\[
 \sum_{n=0}^{\infty}
       \left\{\frac{x_n}{x_n^2+t}-\frac1{x_n}\right\}
       =\psi(a)-\Re\psi(a+i\sqrt t).
\]
Substitution proves \eqref{sp:critical-ordinary-sum}.
The asserted $O(x_n^{-3})$ remainder also follows immediately
from \eqref{sp:connected-asymptotic}.
\end{proof}

\subsection{Polynomial weights and the inherited local correction}
\label{sp:weighted-bridge}

Let $W(x)=\sum_{d=0}^{D}w_dx^d$ be a polynomial.  For functions
with the relevant meromorphic continuation, use the fixed
regulator
\begin{equation}
 \mathcal R_W[f]
 :=\operatorname{FP}_{\varepsilon=0}
     \sum_{n=0}^{\infty}W(x_n)x_n^{-2\varepsilon}f(x_n),
 \qquad
 Q_{W,a}(t):=\mathcal R_W\left[x\longmapsto\frac{1}{x^2+t}\right],
 \label{sp:fixed-regulator}
\end{equation}
where $\operatorname{FP}$ means the constant Laurent coefficient.
The weight is already included in $\mathcal R_W$.

\begin{proposition}[The compact transform with a polynomial weight]
\label{sp:weighted-compact}
For the connected product in \eqref{sp:connected-product},
\begin{equation}
 \mathcal R_W[x\longmapsto G_{\boldsymbol\nu,\beta}(x^2)]
 =\int_\beta^M\rho_{\boldsymbol\nu,\beta}(t)Q_{W,a}(t)\,dt.
 \label{sp:weighted-compact-transform}
\end{equation}
The integral on the right is ordinary and absolutely convergent.
The kernels for monomials have the explicit expressions, for
integers $p\geq0$,
\begin{align}
 Q_{x^{2p},a}(t)
 &=\sum_{j=0}^{p-1}(-t)^j\zeta(2j-2p+2,a)
                        +(-t)^pQ_a(t),
 \label{sp:even-weight-kernel}\\
 Q_{x^{2p+1},a}(t)
 &=\sum_{j=0}^{p-1}(-t)^j\zeta(2j-2p+1,a)
                        -(-t)^p\Re\psi(a+i\sqrt t).
 \label{sp:odd-weight-kernel}
\end{align}
Both sums are empty for $p=0$, and the kernels have their
continuous values at $t=0$.
\end{proposition}

\begin{proof}
For $\Re\varepsilon>(D-1)/2$, the compact representation
\eqref{sp:compact-stieltjes} may be summed with the weight
$W(x_n)x_n^{-2\varepsilon}$ by absolute Fubini.
To continue this identity toward $\varepsilon=0$, choose $K$
sufficiently large and use the exact finite expansion
\[
 \frac{1}{x^2+t}
 =\sum_{j=0}^{K-1}(-t)^j x^{-2j-2}
       +\frac{(-t)^K}{x^{2K}(x^2+t)}.
\]
After summation the finite sum consists of terms
$w_d(-t)^j\zeta(2\varepsilon+2j+2-d,a)$.
The remainder is normally convergent for $\varepsilon$ in a
neighborhood of zero, uniformly on $\beta\leq t\leq M$.
Indeed, its summands have size at most a constant times
$x_n^{D-2\Re\varepsilon-2K-2}$ there.
The possible polar coefficients are polynomials in $t$, and
the holomorphic remainder has a uniform integrable majorant
against $|\rho_{\boldsymbol\nu,\beta}(t)|\,dt$.
Laurent coefficient extraction may therefore be interchanged
with the compact integral, proving
\eqref{sp:weighted-compact-transform}.  This argument uses
meromorphic continuation with uniform control, rather than a
pointwise limit of a divergent unregulated sum.

Polynomial division proves \eqref{sp:even-weight-kernel} from
$Q_{1,a}=Q_a$.  For the odd case, first observe that
\[
 \sum_{n=0}^{\infty}x_n^{-2\varepsilon}
                \frac{x_n}{x_n^2+t}
 =\zeta(1+2\varepsilon,a)
  +\sum_{n=0}^{\infty}x_n^{-2\varepsilon}
        \left\{\frac{x_n}{x_n^2+t}-\frac1{x_n}\right\}.
\]
The second series is normal near zero.  Since the constant
Laurent coefficient of $\zeta(1+2\varepsilon,a)$ is $-\psi(a)$,
the digamma difference used in
Theorem~\ref{sp:critical-weight} gives
$Q_{x,a}(t)=-\Re\psi(a+i\sqrt t)$.
Polynomial division now proves \eqref{sp:odd-weight-kernel}.
\end{proof}

\begin{remark}[Scope and normalization of the local term]
\label{sp:inherited-local-correction}
The ordinary scalar formula \eqref{sp:master-formula} and the
compact transform fit the local-correction calculation already
present in the arithmetic spectral jets continuation
\cite{ProveItArithmeticJets}.  We record its normalization here
as inherited input.  Define
\[
 J_{W,N,r}:=\left.\partial_s^r
       \sum_{n=0}^{\infty}W(x_n)
                     \prod_{j=1}^{N}(x_n^2+\mu_j)^{-s}
                    \right|_{s=0},\qquad
 L(u):=\sum_{j=1}^{N}\log(1+\mu_j u),
\]
with the spectral sum continued before taking its derivative.
For $r\geq1$, the existing correction formula is
\begin{equation}
 \begin{split}
 J_{W,N,r}
 &-\mathcal R_W\left[
        x\longmapsto\left(-\sum_{j=1}^{N}\log(x^2+\mu_j)\right)^r
                   \right]\\
 &=\frac{(-1)^{r+1}}{2N(r+1)}
     \sum_{k\geq r+1}w_{2k-1}[u^k]L(u)^{r+1}.
 \end{split}
 \label{sp:inherited-correction}
\end{equation}
The sum is finite because $W$ is a polynomial.  Its normalization
can be read from the resonant factor
$[u^k]e^{-sL(u)}\zeta(1+2Ns,a)$: the pole contributes
$r![s^{r+1}u^k]e^{-sL(u)}/(2N)$ to the $r$th spectral jet.
This equals the corresponding summand in
\eqref{sp:inherited-correction}.  The fixed regulator has no
such additional constant coefficient from its pole.

For $W=1$ the correction vanishes, consistent with the direct
normal-convergence proof of Theorem~\ref{sp:master}.
More generally, it vanishes for every even polynomial weight.
These are specializations of the inherited formula, not a new
general regularized-trace principle; all-coefficient Laurent
expansions of regularized traces belong to the established
framework represented by Paycha--Scott \cite{PaychaScott}.
Throughout such comparisons, interaction among masses means
replacement of labeled masses at fixed $N$, as in
Remark~\ref{sp:fixed-N-differences}.
\end{remark}


% END sections/04_spectral.tex


% BEGIN sections/04b_negative_jets.tex
\section{Even polynomial weights and jets at negative integers}
\label{ng:section}

The scalar resolvent construction extends to every polynomial in
$x_n^2$, after subtracting a finite asymptotic expansion before
forming the cut density.  Consequently, derivatives at every
negative integer spectral point admit one ordinary digamma
integral.  This is an integral representation in a specified
function class; it does not assert a finite evaluation in
Hurwitz zeta constants or a recurrence eliminating integral
kernels.

Fix $a>0$, $x_n=n+a$, an integer $N\geq1$, and positive masses
$\mu_1,\ldots,\mu_N$.  Let
\[
 P(z)=\sum_{d=0}^{D}p_dz^d\in\mathbb R[z],\qquad
 L(u)=\sum_{j=1}^{N}\log(1+\mu_j u),\qquad
 b_{k,j}:=[u^k]L(u)^j.
\]
Here $L$ is its analytic germ at zero, so
$b_{k,0}=\boldsymbol1_{\{k=0\}}$ and $b_{k,j}=0$ for $k<j$.
We also use
\[
 c_k(s):=[u^k]e^{-sL(u)}
   =\sum_{j=0}^{k}\frac{(-s)^j}{j!}b_{k,j}.
\]
Thus each $c_k$ is a polynomial in $s$, with $c_0=1$ and
$c_k(0)=0$ for $k>0$.  The logarithm $\Log$ below is principal
on $\mathbb C\setminus(-\infty,0]$, and
$\zeta^{(q)}(v,a)=\partial_v^q\zeta(v,a)$.


\subsection{Finite subtraction and holomorphy at zero}
\label{ng:continuation-subsection}

Initially, for $\Re s>(2D+1)/(2N)$, define
\begin{equation}
 Z_{P,N}(s;a,\boldsymbol\mu)
 :=\sum_{n=0}^{\infty}P(x_n^2)
                 \prod_{j=1}^{N}(x_n^2+\mu_j)^{-s}.
 \label{ng:weighted-definition}
\end{equation}
If $P=0$, all statements below are immediate; otherwise $D$
is its degree.

\begin{lemma}[A finite normally convergent continuation]
\label{ng:finite-continuation}
The weighted spectral function has the meromorphic continuation
\begin{align}
 Z_{P,N}(s)
 &=\sum_{d=0}^{D}p_d\sum_{k=0}^{d}
       c_k(s)\zeta(2Ns+2k-2d,a)+R_{P,N}(s),
 \label{ng:continuation}\\
 R_{P,N}(s)
 &:=\sum_{n=0}^{\infty}\sum_{d=0}^{D}p_d x_n^{2d-2Ns}
       \left\{e^{-sL(x_n^{-2})}
                 -\sum_{k=0}^{d}c_k(s)x_n^{-2k}\right\}.
 \label{ng:normal-remainder}
\end{align}
For positive $x_n^{-2}$, the expression $L(x_n^{-2})$ in this
formula means $\sum_j\log(1+\mu_j/x_n^2)$, whether or not it
lies inside the Taylor disk at zero.
The remainder series is normally convergent on
$\Re s>-1/(2N)$.  Every finite zeta term is holomorphic near
zero, so $Z_{P,N}$ is holomorphic there.  In particular,
\begin{equation}
 Z_{P,N}(0)=\sum_{d=0}^{D}p_d\zeta(-2d,a).
 \label{ng:zeroth-jet}
\end{equation}
\end{lemma}

\begin{proof}
For a fixed $d$, Taylor's theorem for the analytic function
$e^{-sL(u)}$ gives
\[
 e^{-sL(u)}-\sum_{k=0}^{d}c_k(s)u^k=O_K(u^{d+1})
 \qquad(u\to0),
\]
uniformly for $s$ in any compact set $K$.
After multiplication by $x_n^{2d-2Ns}$ and substitution
$u=x_n^{-2}$, this is bounded in absolute value by a constant
times $x_n^{-2N\inf_K\Re s-2}$.
This tail is summable if $K\subset\{\Re s>-1/(2N)\}$.
The finitely many small $n$ contribute entire functions of $s$.
Thus \eqref{ng:normal-remainder} is normally convergent in the
asserted half-plane.  Adding and subtracting the indicated
finite powers proves \eqref{ng:continuation} in the initial
half-plane and hence supplies the continuation.
The finite zeta terms can have a pole only when
$2Ns+2k-2d=1$, or
$s=(2(d-k)+1)/(2N)>0$; none has a pole at zero.
Finally $c_k(0)=\boldsymbol1_{\{k=0\}}$ makes the entire
remainder vanish at zero and proves \eqref{ng:zeroth-jet}.
\end{proof}

For integers $r\geq1$, write
$J_{P,N,r}:=\partial_s^rZ_{P,N}(0)$ and define the finite
subtraction
\begin{equation}
 \mathcal A_{P,N,r}(z)
 :=\sum_{d=0}^{D}p_d\sum_{k=0}^{d}\sum_{j=0}^{r}
       \binom{r}{j}N^{r-j}b_{k,j}
                     z^{d-k}(\Log z)^{r-j}.
 \label{ng:asymptotic-subtraction}
\end{equation}
Let
\begin{equation}
 F_{P,N,r}(z)
 :=P(z)\left\{\sum_{j=1}^{N}\Log(z+\mu_j)\right\}^{r}
                  -\mathcal A_{P,N,r}(z),
 \label{ng:subtracted-function}
\end{equation}
and set
\begin{equation}
 B_{P,N,r}
 :=\sum_{d=0}^{D}p_d\sum_{k=0}^{d}\sum_{j=0}^{r}
       (-1)^j\binom{r}{j}(2N)^{r-j}b_{k,j}
                    \zeta^{(r-j)}(-2(d-k),a).
 \label{ng:finite-baseline}
\end{equation}
Terms with $j>k$ vanish automatically.

\begin{proposition}[An ordinary remainder series for every even weight]
\label{ng:ordinary-series}
For every integer $r\geq1$,
\begin{equation}
 J_{P,N,r}
 =B_{P,N,r}+(-1)^r\sum_{n=0}^{\infty}F_{P,N,r}(x_n^2).
 \label{ng:jet-series}
\end{equation}
The series is absolutely convergent.  More precisely, uniformly
at infinity on the slit plane,
\begin{equation}
 F_{P,N,r}(z)
 =O\left(\frac{1+\log^{r-1}|z|}{|z|}\right).
 \label{ng:large-z-bound}
\end{equation}
\end{proposition}

\begin{proof}
Normal convergence in Lemma~\ref{ng:finite-continuation}
permits differentiation at zero.  Since
$c_k^{(j)}(0)=(-1)^j b_{k,j}$, the derivative of its finite
zeta sum is exactly \eqref{ng:finite-baseline}.
In the remainder, the derivative of the original summand is
\[
 (-1)^r x_n^{2d}
       \left\{\sum_{j=1}^{N}\log(x_n^2+\mu_j)\right\}^{r},
\]
while the derivative of
$x_n^{2d-2k-2Ns}c_k(s)$ is
\[
 (-1)^r x_n^{2d-2k}\sum_{j=0}^{r}
             \binom{r}{j}N^{r-j}b_{k,j}(\log x_n^2)^{r-j}.
\]
This proves \eqref{ng:jet-series}, including its signs and
powers of $2N$.

For the precise decay, take $|z|$ larger than the largest mass.
Then, with branches continued from the positive real axis,
\[
 \sum_j\Log(z+\mu_j)=N\Log z+L(1/z).
\]
For $j\geq1$, introduce the Taylor remainder
\[
 T_{d,j}(u):=L(u)^j-\sum_{k=0}^{d}b_{k,j}u^k
                    =O(u^{d+1}).
\]
The $j=0$ contribution is removed exactly by
\eqref{ng:asymptotic-subtraction}, leaving the exact identity
\begin{equation}
 F_{P,N,r}(z)
 =\sum_{d=0}^{D}p_d\sum_{j=1}^{r}
       \binom{r}{j}N^{r-j}z^d(\Log z)^{r-j}T_{d,j}(1/z)
 \label{ng:tail-decomposition}
\end{equation}
in this exterior domain.  Each term is
$O(|z|^{-1}(1+\log^{r-1}|z|))$.
This proves \eqref{ng:large-z-bound}, and hence the ordinary
series is dominated by a constant times
$x_n^{-2}(1+\log^{r-1}x_n)$ for large $n$.
\end{proof}


\subsection{The one-integral theorem for every even weight}
\label{ng:integral-subsection}

For $t>0$ away from the masses, put
\begin{equation}
 A(t):=\sum_{j=1}^{N}\log|\mu_j-t|,\qquad
 k(t):=\#\{j:\mu_j<t\}.
 \label{ng:cut-data}
\end{equation}
Repeated masses are counted with their multiplicity.
The upper-cut value of the finite subtraction is explicitly
\begin{equation}
 \mathcal A_{P,N,r}^{+}(t)
 :=\sum_{d=0}^{D}p_d\sum_{k=0}^{d}\sum_{j=0}^{r}
       \binom{r}{j}N^{r-j}b_{k,j}
                 (-t)^{d-k}(\log t+i\pi)^{r-j}.
 \label{ng:upper-subtraction}
\end{equation}
All powers $(-t)^{d-k}$ in this expression are real integer
powers.  Define
\begin{equation}
 \rho_{P,N,r}(t)
 :=-\frac1\pi\Im\left\{
       P(-t)[A(t)+i\pi k(t)]^r
                    -\mathcal A_{P,N,r}^{+}(t)\right\}.
 \label{ng:density}
\end{equation}
The values of the density at the mass points are immaterial.

\begin{theorem}[Even polynomial weights in one ordinary integral]
\label{ng:master}
For every $P\in\mathbb R[z]$, every $N\geq1$, every $r\geq1$,
$a>0$, and positive masses,
\begin{equation}
 \boxed{\displaystyle
 J_{P,N,r}=B_{P,N,r}
          +(-1)^r\int_0^{\infty}\rho_{P,N,r}(t)Q_a(t)\,dt,}
 \label{ng:master-formula}
\end{equation}
where the finite baseline is \eqref{ng:finite-baseline} and
\begin{equation}
 Q_a(t)=\sum_{n=0}^{\infty}\frac1{(n+a)^2+t}
       =\frac{\psi(a+i\sqrt t)-\psi(a-i\sqrt t)}{2i\sqrt t},
 \qquad Q_a(0)=\psi'(a).
 \label{ng:digamma-kernel}
\end{equation}
The integral is absolutely convergent.  Splitting it at the
distinct masses produces ordinary improper integrals, with no
finite-part or principal-value prescription.

In addition, the subtracted function itself satisfies
\begin{equation}
 F_{P,N,r}(z)=\int_0^{\infty}
                         \frac{\rho_{P,N,r}(t)}{z+t}\,dt,
 \qquad z\in\mathbb C\setminus(-\infty,0].
 \label{ng:stieltjes-representation}
\end{equation}
For $r=1$, the density simplifies to the compactly supported
expression
\begin{equation}
 \rho_{P,N,1}(t)=P(-t)\{N-k(t)\}
       =P(-t)\sum_{j=1}^{N}\boldsymbol1_{(0,\mu_j)}(t).
 \label{ng:first-density}
\end{equation}
\end{theorem}

\begin{proof}
The function in \eqref{ng:subtracted-function} is holomorphic
on the slit plane.  All coefficients and masses are real, so
it satisfies Schwarz reflection.  On the upper cut,
\[
 \sum_j\Log(\mu_j-t+i0)=A(t)+i\pi k(t),\qquad
 \Log(-t+i0)=\log t+i\pi.
\]
Thus its upper boundary value has imaginary part
$-\pi\rho_{P,N,r}(t)$, with exactly the sign in
\eqref{ng:density}.

At zero the masses are positive and the first term in
\eqref{ng:subtracted-function} stays bounded.  The finite
subtraction contains only nonnegative powers of $z$, and hence
\[
 F_{P,N,r}(z)=O(1+|\Log z|^r)\qquad(z\to0).
\]
At a nonzero cut endpoint $-\mu_j$, the singularity is at worst
a power of a logarithm.  Therefore the Cauchy integrals over
small endpoint circles vanish, since
$\varepsilon(1+|\log\varepsilon|^r)\to0$.
The large-circle contribution vanishes by
\eqref{ng:large-z-bound}.  Collapsing a Cauchy contour onto the
two banks of the cut gives \eqref{ng:stieltjes-representation}.

For completeness, the following bounds establish absolute
convergence and justify summing that representation.  Near zero,
$k(t)=0$, so the first expression inside the imaginary part in
\eqref{ng:density} is real.  Taking the imaginary part of the
finite subtraction lowers each positive logarithmic degree
by one.  Consequently,
\begin{equation}
 \rho_{P,N,r}(t)=O(1+|\log t|^{r-1})\qquad(t\downarrow0).
 \label{ng:zero-bound}
\end{equation}
At any mass point it has at worst an integrable logarithmic
singularity; in particular a local bound is a constant times
$1+|\log|t-\mu_j||^{r-1}$.
To control infinity without estimating separately the large
terms of \eqref{ng:density}, use the exact tail expression
\eqref{ng:tail-decomposition}.
For $t>\max_j\mu_j$, all numbers $T_{d,j}(-1/t)$ are real.
Its only imaginary parts come from
$(\log t+i\pi)^{r-j}$.  The term with $r-j=0$ is real; the
remaining terms have imaginary parts of logarithmic degree at
most $r-j-1\leq r-2$.  Therefore, for $r\geq2$,
\begin{equation}
 \rho_{P,N,r}(t)
       =O\left(\frac{1+\log^{r-2}t}{t}\right)
                          \qquad(t\to\infty).
 \label{ng:infinity-bound}
\end{equation}
For $r=1$ every exterior tail term is real, so the density
vanishes beyond the largest mass.  More explicitly,
\eqref{ng:asymptotic-subtraction} at $r=1$ is
$NP(z)\Log z$ plus a polynomial with real coefficients.
Its upper-cut imaginary part is $\pi NP(-t)$; comparison with
the first term in \eqref{ng:density} proves
\eqref{ng:first-density} on every interval of the cut.

The convergent digamma difference series proves
\eqref{ng:digamma-kernel}.  Its series expression shows that
$Q_a$ is positive and bounded near zero.  Integral comparison
gives $Q_a(t)=\pi/(2\sqrt t)+O(t^{-1})$ at infinity.
Combining these facts with
\eqref{ng:zero-bound}--\eqref{ng:infinity-bound}, or with compact
support for $r=1$, yields
\[
 \int_0^{\infty}|\rho_{P,N,r}(t)|Q_a(t)\,dt<\infty.
\]
Hence absolute Fubini gives
\[
 \sum_{n=0}^{\infty}F_{P,N,r}(x_n^2)
  =\int_0^{\infty}\rho_{P,N,r}(t)
                        \sum_{n=0}^{\infty}\frac{dt}{x_n^2+t}
  =\int_0^{\infty}\rho_{P,N,r}(t)Q_a(t)\,dt.
\]
Together with \eqref{ng:jet-series}, this is
\eqref{ng:master-formula}.
\end{proof}

\begin{remark}[Collisions and the role of real coefficients]
\label{ng:collisions}
No distinctness assumption is used: coincident masses simply
increase the multiplicities in $A(t)$ and $k(t)$, while the
endpoint singularities remain integrable.  The coefficients
$p_d$ need not be positive; the density can be signed.
Their reality is used in Schwarz reflection and in the improved
imaginary-part estimate at infinity.  The theorem is stated for
real polynomial weights with precisely these properties.
For $P=1$ it reduces to the scalar master, since
$\mathcal A_{1,N,r}=N^r\Log^r z$ and
$B_{1,N,r}=(2N)^r\zeta^{(r)}(0,a)$.
\end{remark}


\subsection{All negative integer spectral points}
\label{ng:negative-subsection}

\begin{corollary}[Jets at every negative integer]
\label{ng:negative-integers}
For an integer $m\geq0$, define the real polynomial
\begin{equation}
 P_m(z):=\prod_{j=1}^{N}(z+\mu_j)^m.
 \label{ng:negative-weight}
\end{equation}
Its coefficients are held fixed when differentiating in the
spectral variable.  The scalar function
\[
 Z_N(v;a,\boldsymbol\mu)
   =\sum_{n=0}^{\infty}\prod_{j=1}^{N}(x_n^2+\mu_j)^{-v}
\]
is holomorphic at $v=-m$, and for every integer $r\geq1$,
\begin{equation}
 \boxed{\displaystyle
 Z_N^{(r)}(-m;a,\boldsymbol\mu)
 =B_{P_m,N,r}
       +(-1)^r\int_0^{\infty}\rho_{P_m,N,r}(t)Q_a(t)\,dt.}
 \label{ng:negative-master}
\end{equation}
The baseline and density are the explicit finite formulas
\eqref{ng:finite-baseline} and \eqref{ng:density}, with
$P=P_m$ and degree $mN$.  The integral is ordinary and absolutely
convergent, including at coincident masses.
Writing $P_m(z)=\sum_{d=0}^{mN}p_{m,d}z^d$, the value itself is
\begin{equation}
 Z_N(-m;a,\boldsymbol\mu)
       =\sum_{d=0}^{mN}p_{m,d}\zeta(-2d,a).
 \label{ng:negative-value}
\end{equation}
More generally, if a fixed even weight is $E(x_n^2)$ with
$E\in\mathbb R[z]$, the same formulas give every derivative at
$-m$ after replacing $P_m$ by $E P_m$.
\end{corollary}

\begin{proof}
In the common half-plane of absolute convergence,
\[
 Z_{P_m,N}(s)
 =\sum_{n=0}^{\infty}\prod_{j=1}^{N}
                      (x_n^2+\mu_j)^{m-s}
 =Z_N(s-m).
\]
Lemma~\ref{ng:finite-continuation} continues the left side to
$\Re s>-1/(2N)$ and proves that it is holomorphic at zero.
After shifting the spectral variable, this gives a continuation
of $Z_N(v)$ through $v=-m$.  Continuations constructed with
different $m$ agree on their overlaps by uniqueness.
Taking the $r$th derivative at $s=0$ and applying
Theorem~\ref{ng:master} proves \eqref{ng:negative-master}.
The value formula follows from \eqref{ng:zeroth-jet}.
With an additional fixed weight, the identical calculation is
$Z_{E P_m,N}(s)=Z_{E,N}(s-m)$, which proves the final assertion.
\end{proof}

\subsection{An explicit cubic jet at the first negative point}
\label{ng:cubic-example-subsection}

Take $N=1$, a single mass $\mu>0$, $m=1$, and $r=3$.
Then $P_1(z)=z+\mu$ and $L(u)=\log(1+\mu u)$.
The only coefficient with $1\leq k\leq1$ needed here is
$b_{1,1}=\mu$.  Directly from
\eqref{ng:asymptotic-subtraction}--\eqref{ng:finite-baseline},
\begin{align}
 \mathcal A_{z+\mu,1,3}(z)
   &=(z+\mu)\Log^3z+3\mu\Log^2z,
 \label{ng:cubic-subtraction}\\
 F_{z+\mu,1,3}(z)
   &=(z+\mu)\Log^3(z+\mu)
          -(z+\mu)\Log^3z-3\mu\Log^2z,
 \label{ng:cubic-remainder}\\
 B_{z+\mu,1,3}
   &=8\zeta^{(3)}(-2,a)+8\mu\zeta^{(3)}(0,a)
                       -12\mu\zeta^{(2)}(0,a).
 \label{ng:cubic-baseline}
\end{align}
The density is the following elementary piecewise expression:
\begin{equation}
 \rho_\mu(t)=
 \begin{cases}
  (\mu-t)(3\log^2t-\pi^2)+6\mu\log t,
        &0<t<\mu,\\[3pt]
  3(\mu-t)\{\log^2t-\log^2(t-\mu)\}+6\mu\log t,
        &t>\mu.
 \end{cases}
 \label{ng:cubic-density}
\end{equation}
It has the common limiting value $6\mu\log\mu$ at $t=\mu$.
Thus the complete ordinary-integral formula is
\begin{equation}
 \begin{split}
 Z_1^{(3)}(-1;a,\mu)
   ={}&8\zeta^{(3)}(-2,a)+8\mu\zeta^{(3)}(0,a)
                          -12\mu\zeta^{(2)}(0,a)\\
      &-\int_0^{\infty}\rho_\mu(t)Q_a(t)\,dt.
 \end{split}
 \label{ng:cubic-negative-jet}
\end{equation}

\begin{proof}
The displayed finite subtraction and baseline follow by
retaining $d=0,1$ and $k\leq d$ in the general formulas.
For $0<t<\mu$, the value $\Log(\mu-t+i0)$ is real, so only
the subtracted terms contribute to the density.  The identities
\[
 \frac1\pi\Im(\log t+i\pi)^3=3\log^2t-\pi^2,
 \qquad
 \frac1\pi\Im(\log t+i\pi)^2=2\log t
\]
give its first branch.  For $t>\mu$, the original logarithm
is $\log(t-\mu)+i\pi$.  Its cubic imaginary part supplies the
second logarithmic square, and the two $\pi^2$ terms cancel.
This gives the second branch of \eqref{ng:cubic-density}.
The limit at $t=\mu$ uses
$(t-\mu)\log^2|t-\mu|\to0$.
Finally $(-1)^3=-1$ in \eqref{ng:negative-master}, which gives
the minus sign before the integral in
\eqref{ng:cubic-negative-jet}.
\end{proof}

\begin{remark}[A tail check and the scope of the extension]
\label{ng:scope}
For fixed $\mu>0$, expansion of the second branch yields
\begin{equation}
 \rho_\mu(t)=\frac{3\mu^2(\log t+1)}{t}
                  +O\left(\frac{1+\log t}{t^2}\right)
                       \qquad(t\to\infty),
 \label{ng:cubic-tail}
\end{equation}
which also checks the cancellation of all apparent growing
terms in \eqref{ng:cubic-density}.
For example, putting $\delta=\log(1-\mu/t)$ rewrites that
branch exactly as
\[
 6\log t\{\mu+(t-\mu)\delta\}+3(t-\mu)\delta^2;
\]
Taylor expansion then gives \eqref{ng:cubic-tail}.

The general construction supplies, at every negative integer,
a finite Hurwitz-derivative baseline and one ordinary integral
against the same digamma resolvent.  The new density is explicit
in finite Taylor coefficients and logarithmic boundary values.
For orders $r\geq2$ it generally has an infinite tail, controlled
by \eqref{ng:infinity-bound}; at order one it has compact support.
This resolves the representation question within the ordinary
digamma-integral function class.  Finite evaluation in a smaller
set of constants, or the separate recurrence for half-integer Bessel-order
derivatives, requires additional identities.
\end{remark}

% END sections/04b_negative_jets.tex


% BEGIN sections/05_density_arithmetic.tex
\section{Mellin moments, harmonic collisions, and polylogarithmic transforms}
\label{da:section}

The compact density has a useful arithmetic structure. This section
identifies it explicitly and shows how the same polynomials connect the
spectral formulas to Gamma ratios, polygamma differences, generalized
harmonic numbers, and multiple polylogarithms. The coefficient identities
and height-one reductions used in this bridge are classical. Their role
here is to give computable consequences of the logarithmic density, not
to assert a new historical origin for those identities.

\subsection{The beta generating function of the density}

Retain the polynomials \(P_r\) of \eqref{sp:coincident-polynomial} and put
\[
 y(t)=\log\frac{1-t}{t},\qquad 0<t<1.
\]
Since \(P_r(y)=-\pi^{-1}\Im(y-i\pi)^r\) for real \(y\), coefficient
expansion gives the entire generating identity
\begin{equation}\label{da:Pgen}
 \sum_{r\ge1}P_r(y)\frac{w^r}{r!}
       =e^{wy}\frac{\sin\pi w}{\pi}.
\end{equation}
For complex \(w\), the right side is the analytic continuation of this
polynomial coefficient identity; no imaginary-part operation is applied
to a complex parameter \(w\).

\begin{theorem}[Every Mellin moment of the logarithmic density]
\label{da:mellin}
For \(\Re s>0\), define the ordinary integral
\[
 M_r(s)=\int_0^1 t^{s-1}P_r(y(t))\dd t,\qquad r\ge1.
\]
For \(w\) in a neighborhood of zero,
\begin{equation}\label{da:mellin-gen}
 \boxed{\displaystyle
  \sum_{r\ge1}M_r(s)\frac{w^r}{r!}
     =\frac{w\Gamma(s-w)}{\Gamma(s+1)\Gamma(1-w)}.}
\end{equation}
More generally, the underlying integral of \eqref{da:Pgen} equals this
Gamma ratio whenever \(\Re(s-w)>0\) and \(\Re(1+w)>0\).
\end{theorem}

\begin{proof}
On a sufficiently small closed disk about \(w=0\), the integrand and all
fixed \(w\)-derivatives are bounded by
\[
 C\,t^{\Re s-1-\eta}(1-t)^{-\eta}
       (1+|\log t|^j+|\log(1-t)|^j),
\]
with \(0<\eta<\min(1,\Re s)\). This is integrable. Thus
coefficient extraction may be interchanged with the integral.
Euler's beta integral and the Gamma reflection relation then give
\begin{align*}
 \int_0^1t^{s-1}e^{wy(t)}\frac{\sin\pi w}{\pi}\dd t
 &=\frac{\sin\pi w}{\pi}
      \frac{\Gamma(s-w)\Gamma(1+w)}{\Gamma(s+1)}\\
 &=\frac{w\Gamma(s-w)}{\Gamma(s+1)\Gamma(1-w)}.
\end{align*}
This proves both the generating identity and its stated domain.
\end{proof}

Let \(Y_n\) denote the complete exponential Bell polynomials, normalized by
\[
 \exp\left(\sum_{j\ge1}a_j\frac{w^j}{j!}\right)
      =\sum_{n\ge0}Y_n(a_1,\ldots,a_n)\frac{w^n}{n!},
 \qquad Y_0=1.
\]
Taking the logarithm of the Gamma ratio in \eqref{da:mellin-gen} gives
the finite polygamma formula
\begin{equation}\label{da:polygamma-moment}
 M_r(s)=\frac r{s}Y_{r-1}(a_1(s),\ldots,a_{r-1}(s)),
 \quad
 a_j(s)=(-1)^j\{\psi^{(j-1)}(s)-\psi^{(j-1)}(1)\}.
\end{equation}
In particular,
\begin{align*}
 M_1(s)&=\frac1s,\\
 M_2(s)&=\frac2s\{\psi(1)-\psi(s)\},\\
 M_3(s)&=\frac3s\bigl\{(\psi(s)-\psi(1))^2+
                                   \psi'(s)-\psi'(1)\bigr\}.
\end{align*}
All \(s\)-derivatives of these identities are ordinary logarithmic Mellin
moments, justified by the same endpoint majorants.

\begin{corollary}[Stirling moments and exact orthogonality]
\label{da:stirling}
Let \(\stir nr\) be the unsigned Stirling number of the first kind, with
\(\stir nr=0\) when \(n<r\). For every \(k\ge0\), \(r\ge1\),
\begin{equation}\label{da:stirling-moment}
 \boxed{\displaystyle
  \int_0^1 t^kP_r\!\left(\log\frac{1-t}{t}\right)\dd t
    =(-1)^{r-1}\frac{r!}{(k+1)!}\stir{k+1}{r}.}
\end{equation}
Thus \(P_r(y(t))\) is orthogonal to every polynomial in \(t\) of degree
less than \(r-1\), with respect to ordinary Lebesgue integration.
\end{corollary}

\begin{proof}
At \(s=k+1\), the right side of \eqref{da:mellin-gen} is
\[
 \frac{w(1-w)_k}{(k+1)!}.
\]
Expanding this finite polynomial using
\(u(u+1)\cdots(u+k)=\sum_r\stir{k+1}{r}u^r\)
gives \eqref{da:stirling-moment} with its sign. The asserted
orthogonality follows from the vanishing Stirling numbers.
\end{proof}

This moment formula also makes the harmonic content explicit:
\[
 \frac{\stir{k+1}{r}}{k!}
 =e_{r-1}\left(1,\frac12,\ldots,\frac1k\right),
\]
where \(e_j\) is the elementary symmetric polynomial, and \(e_0=1\).
Equivalently, \eqref{da:polygamma-moment} at \(s=k+1\) has
\(a_j=-(j-1)!H_k^{(j)}\). These finite formulas make exact moment
calculation possible without numerical quadrature.

\subsection{A transform from ordinary to height-one multiple polylogarithms}

For \(|z|<1\) and \(s\in\mathbb C\), use
\[
 \Li_s(z)=\sum_{n\ge1}\frac{z^n}{n^s}.
\]
The height-one multiple polylogarithm needed here can be defined directly
by the locally normally convergent series
\begin{equation}\label{da:heightone-definition}
 \Li_{s+1,\one{r-1}}(z)
   =\sum_{n\ge1}\frac{z^n}{n^{s+1}}H_{n-1}(\one{r-1}),
\end{equation}
where
\[
 H_m(\one d)=
   \sum_{m\ge n_1>\cdots>n_d\ge1}\frac1{n_1\cdots n_d},
 \qquad H_m(\one0)=1.
\]
Thus the nesting convention is specified rather than inferred from
notation.

\begin{theorem}[The logarithmic-density polylogarithm transform]
\label{da:polylog}
For every \(r\ge1\), every \(s\in\mathbb C\), and \(|z|<1\),
\begin{equation}\label{da:polylog-transform}
 \boxed{\displaystyle
 \int_0^1 P_r\!\left(\log\frac{1-t}{t}\right)
                \frac{\Li_s(zt)}t\dd t
   =(-1)^{r-1}r!\,\Li_{s+1,\one{r-1}}(z).}
\end{equation}
The integral and series are locally normal in \(s\). Consequently every
order derivative satisfies
\begin{equation}\label{da:polylog-order-jets}
 \int_0^1P_r(y(t))\frac{\partial_s^j\Li_s(zt)}t\dd t
  =(-1)^{r-1}r!\,
       \partial_s^j\Li_{s+1,\one{r-1}}(z),\qquad j\ge0.
\end{equation}
For integer \(s\ge1\), the ordinary endpoint value at \(z=1\) is
\begin{equation}\label{da:mzv-transform}
 \int_0^1P_r(y(t))\frac{\Li_s(t)}t\dd t
   =(-1)^{r-1}r!\,\zeta(s+1,\one{r-1}).
\end{equation}
\end{theorem}

\begin{proof}
For \(|z|\le R<1\) and \(s\) in a compact set, the polylogarithm
series and any fixed order derivative are dominated by
\(\sum_n R^n n^A(1+\log^j n)\). Near zero,
\(\Li_s(zt)/t\) remains bounded uniformly on these sets; near one
the only additional factor is an integrable polynomial in
\(\log(1-t)\). Absolute termwise integration and
\eqref{da:stirling-moment} give
\[
 (-1)^{r-1}r!\sum_{n\ge1}
       \frac{z^n}{n^s n!}\stir nr.
\]
The finite identity
\(\stir nr/(n-1)!=H_{n-1}(\one{r-1})\) turns this into
\eqref{da:heightone-definition}. The same domination justifies the
order derivatives.

For real \(0\le z<1\) and integer \(s\ge1\),
\(|\Li_s(zt)|\le\Li_s(t)\). At \(s=1\) this grows at most as
\(|\log(1-t)|\), and at \(s>1\) it is bounded near one.
Dominated convergence gives the endpoint integral. The positive series
in \eqref{da:heightone-definition} converges at one because its first
index is at least two, proving \eqref{da:mzv-transform}.
\end{proof}

\begin{corollary}[An all-order ordinary zeta integral]
\label{da:zeta-integral}
For every integer \(r\ge1\),
\begin{equation}\label{da:zeta-integral-formula}
 \int_0^1P_r\!\left(\log\frac{1-t}{t}\right)
       \frac{-\log(1-t)}t\dd t
     =(-1)^{r-1}r!\,\zeta(r+1).
\end{equation}
For example,
\begin{equation}\label{da:pi-four}
 \boxed{\displaystyle
 \int_0^1\left(3\log^2\frac{1-t}{t}-\pi^2\right)
             \frac{-\log(1-t)}t\dd t=\frac{\pi^4}{15}.}
\end{equation}
\end{corollary}

\begin{proof}
The all-one iterated integral is
\(\Li_{\one r}(z)=(-\log(1-z))^r/r!\), by induction using its
derivative and its zero value at \(z=0\). Integration once against
\(dz/z\), followed by \(u=-\log(1-z)\), gives
\[
 \zeta(2,\one{r-1})
  =\frac1{r!}\int_0^1\frac{(-\log(1-z))^r}{z}\dd z
  =\frac1{r!}\int_0^\infty\frac{u^r}{e^u-1}\dd u
  =\zeta(r+1).
\]
The last equality follows by expanding \(1/(e^u-1)\) into its
positive exponential series and integrating. Apply
\eqref{da:mzv-transform} at \(s=1\). The case \(r=3\), together with
\(\zeta(4)=\pi^4/90\), gives \eqref{da:pi-four}.
\end{proof}

\subsection{Every confluent logarithmic coefficient}

For \(a>0\), \(\beta\ge0\), define the convergent resolvents
\[
 R_k(a;\beta)=\sum_{n\ge0}\frac1{((n+a)^2+\beta)^k}
   =\frac{(-1)^{k-1}}{(k-1)!}\partial_\beta^{\,k-1}Q_a(\beta),
 \qquad k\ge1.
\]
For \(r\) independent displacements satisfying
\(|\delta_j|<a^2+\beta\), the ordinary product sum has the convergent
collision expansion
\begin{align}
 &\sum_{n\ge0}\prod_{j=1}^r
     \log\left(1+\frac{\delta_j}{(n+a)^2+\beta}\right)\notag\\
 &\quad=
 \sum_{k_1,\ldots,k_r\ge1}
   \frac{(-1)^{k_1+\cdots+k_r-r}
             \delta_1^{k_1}\cdots\delta_r^{k_r}}
        {k_1\cdots k_r}\,
       R_{k_1+\cdots+k_r}(a;\beta).
 \label{da:independent-collision}
\end{align}
All branches are the ones obtained from the convergent logarithm series
at zero. On smaller closed polydisks, the series is absolutely and normally
convergent: the absolute logarithm series is bounded by a constant times
\(((n+a)^2+\beta)^{-1}\) in each factor, and the product is summable.
This proves the interchange and permits every fixed parameter derivative.

When all displacements equal \(\delta\), the classical logarithm-power
coefficients give a particularly short all-order formula:
\begin{equation}\label{da:confluent-stirling}
 \boxed{\displaystyle
 \sum_{n\ge0}\log^r\left(1+\frac{\delta}{(n+a)^2+\beta}\right)
   =r!\sum_{k\ge r}\frac{(-1)^{k-r}}{k!}\stir kr
                          \delta^kR_k(a;\beta).}
\end{equation}
For \(r=2\), this recovers the harmonic coefficient in the inherited
second-jet collision expansion \cite{ProveItHigherDefects}. For \(r=3\),
\begin{align}
 &\sum_{n\ge0}\log^3\left(1+\frac{\delta}{(n+a)^2+\beta}\right)\notag\\
 &\qquad=
 3\sum_{k\ge3}\frac{(-1)^{k-3}\delta^k}{k}
       \{H_{k-1}^2-H_{k-1}^{(2)}\}R_k(a;\beta).
 \label{da:cubic-harmonic}
\end{align}
At \(\beta=0\), \(R_k(a;0)=\zeta(2k,a)\); at an integer shift
\(a=M\), it is \(\zeta(2k)-H_{M-1}^{(2k)}\).
Thus the compact density, the convergent mass expansion, and the
polygamma expression for \(R_k\) are three independently useful
representations of the same all-order logarithmic sum.

\subsection{Half-lattices and bilateral polylogarithmic kernels}

Digamma reflection gives the particularly simple half-lattice kernel
\begin{equation}\label{da:half-lattice}
 Q_{1/2}(t)=\frac{\pi}{2\sqrt t}\tanh(\pi\sqrt t),
 \qquad Q_{1/2}(0)=\frac{\pi^2}{2}.
\end{equation}
For example the cubic repeated-mass identity becomes
\begin{align}
 &\sum_{n\ge0}\log^3
    \frac{(n+\frac12)^2+\nu}{(n+\frac12)^2+\beta}\notag\\
 &\quad=\frac{\pi}{2}\int_\beta^\nu
    \left\{3\log^2\frac{\nu-t}{t-\beta}-\pi^2\right\}
                  \frac{\tanh(\pi\sqrt t)}{\sqrt t}\dd t,
       \qquad \nu>\beta\ge0 .
 \label{da:half-cubic}
\end{align}
The apparent endpoint at \(t=0\) is removable in the hyperbolic kernel,
so this remains an ordinary integral when \(\beta=0\).

For \(0<a<1\), define the bilateral resolvent
\[
 Q_a^{\mathrm{bil}}(t)=Q_a(t)+Q_{1-a}(t).
\]
Reflection and the convergent geometric Fourier series give
\begin{align}
 Q_a^{\mathrm{bil}}(t)
 &=\frac{\pi}{\sqrt t}
     \frac{\sinh(2\pi\sqrt t)}
          {\cosh(2\pi\sqrt t)-\cos(2\pi a)}\notag\\
 &=\frac{\pi}{\sqrt t}
       \left\{1+2\Re\Li_0(e^{-2\pi\sqrt t+2\pi ia})\right\},
             \qquad t>0. \label{da:bilateral-kernel}
\end{align}
Here \(\Li_0(z)=z/(1-z)\), so the kernel is elementary.
The bilateral full spectral jet
\(J_{N,r}^{\mathrm{bil}}(a)=J_{N,r}(a)+J_{N,r}(1-a)\) therefore is
\begin{align}
 J_{N,r}^{\mathrm{bil}}(a;\boldsymbol\mu)
 ={}&(2N)^r\{\zeta^{(r)}(0,a)+\zeta^{(r)}(0,1-a)\}\notag\\
 &+(-1)^r\int_0^\infty
          \rho_{N,r}(t)Q_a^{\mathrm{bil}}(t)\dd t.
 \label{da:bilateral-jet}
\end{align}
This is a finite polylogarithmic-integrand description at every order,
while retaining the displayed mass-zero Hurwitz derivatives. It does not
assert a finite evaluation of the remaining mass integral.

There is also an explicit anchored shift primitive of the bilateral
kernel. For \(0<a,a_*<1\), put \(q_a(t)=e^{-2\pi\sqrt t+2\pi ia}\). Then
\begin{equation}\label{da:bilateral-primitive}
 \int_{a_*}^{a}Q_u^{\mathrm{bil}}(t)\dd u
   =\frac{\pi(a-a_*)+
          \Im\{\Li_1(q_a(t))-\Li_1(q_{a_*}(t))\}}{\sqrt t}.
\end{equation}
Indeed \(\partial_a\Li_1(q_a)=2\pi i\Li_0(q_a)\).
The continuous value at \(t=0\) is
\(\pi\{\cot(\pi a_*)-\cot(\pi a)\}\).
Integration of \eqref{da:bilateral-jet} or of its compact connected
version against the shift therefore gives anchored ordinary
polylogarithmic integrals. The required interchanges follow from the
locally uniform bounds already proved for \(Q_a\) and its reflected
partner.

% END sections/05_density_arithmetic.tex


% BEGIN sections/06_audit.tex
\section{Corrections and reconciliation of the existing material}
\label{au:section}

The audit distinguishes a definite transcription error from research-status
updates and from questions that remain unresolved. The supplied patch makes
only the definite correction. The integration note proposes the status
changes separately.

\subsection{A reciprocal-Gamma coefficient contains a transcription error}

In the canonical chapter
\src{07-digamma-power-moments.tex}, equation
\src{dgp:eq:a5}, the final term of the displayed numerator is
\verb|+24\zeta\alpha_5|.
It should be \(24\zeta_5\), where the chapter uses
\(\zeta_j=\zeta(j)\). The corrected coefficient is
\begin{equation}\label{au:alpha5}
 \boxed{\displaystyle
 \alpha_5=
 \frac{\gamma^5-10\gamma^3\zeta(2)+20\gamma^2\zeta(3)
       +15\gamma\zeta(2)^2-30\gamma\zeta(4)
       -20\zeta(2)\zeta(3)+24\zeta(5)}{120}.}
\end{equation}
Indeed, the logarithm of the reciprocal Gamma function is
\[
 \log\frac1{\Gamma(1+u)}
   =\gamma u-\frac{\zeta(2)}2u^2+\frac{\zeta(3)}3u^3
          -\frac{\zeta(4)}4u^4+\frac{\zeta(5)}5u^5+O(u^6).
\]
Exponentiating and extracting \(u^5\) gives
\eqref{au:alpha5}. All other terms in the source's numerator agree with
this calculation. This is a transcription correction, not a new
reciprocal-Gamma identity. The companion patch was checked against the
pinned source with \src{git apply --check}; it was not applied to the
repository.

\subsection{The diagonal negative-order target was already completed}

Research question R7 in \emph{Tail Kernels and Boundary Products}
\cite{ProveItTailKernels} asks for the continued alternating-square law
at nonpositive orders, every order derivative, and the distinction between
ordinary summation and cutoff continuation. The incoming
\emph{Bilinear Tails and Determinants} report already supplies precisely
that diagonal result in
\src{sections/04_cutoff_jets.tex}, including its finite endpoint,
even-power expansion, contacts, translated cutoffs, and convergent
subtractions \cite{ProveItBilinear}.

The appropriate update is to cross-reference that proof and retain the
genuinely broader independent-order question. The latter was explicitly
left open in the same report and is completed by
Theorem~\ref{ct:contacts} here. The diagonal result is also an exact
computational compatibility check on the present bivariate coefficient
algorithm.

\subsection{The cubic Stieltjes target has a preserved general formula}

R9 of \emph{Tail Kernels and Boundary Products} asks for an explicit
holomorphic generating germ for cubic Stieltjes products with independent
indices and positive total argument-derivative order. Its stated first
target is
\[
 \FP\int_0^1\gamma_1(x)^2\gamma_0'(x)\dd x.
\]
The preserved \emph{Coincident Directional Zeta} continuation already
contains a general local completion for any number of factors and a finite
cubic Tornheim--zeta coordinate formula
\cite{ProveItCubic}. Its relevant source labels are
\src{cub:general-completion},
\src{cub:all-moments}, and \src{cub:global}.
Thus the existence of an explicit finite generating formula should be
marked as source-proved, subject to the project's normal integration and
review process. The further question of reducing selected retained
Tornheim order derivatives to a smaller arithmetic alphabet remains
meaningful. The supplied source does not prove general depth-one closure
or arithmetic independence.

The local mechanism was checked independently as follows. Write
\[
 H_{u,r}(x)=\partial_x^r\left(\zeta(1+u,x)-\frac1u\right).
\]
Its singular term is
\[
 (-1)^r(1+u)_r x^{-1-u-r},
\]
and its remainder is jointly holomorphic near \(u=x=0\). In a product
of three factors, choosing the singular term in a nonempty proper subset
\(S\) produces a resonance exactly at Taylor degree
\begin{equation}\label{au:resonant-degree}
  \sum_{i\in S}(1+r_i)-1
\end{equation}
of the remaining regular factors. The corresponding monomial is
\(A_S(\boldsymbol u)x^{-1-\sum_{i\in S}u_i}\). Its analytically
continued integral is
\(-A_S(\boldsymbol u)/\sum_{i\in S}u_i\), while all its order derivatives
at the resonance have fixed-coordinate finite part zero.

For \((r_1,r_2,r_3)=(0,0,1)\), the three singleton Taylor degrees are
\(0,0,1\); the three pair degrees are \(1,2,2\). These are all six
proper nonempty subsets. The full triple has no remaining regular factor
and has baseline exponent \(-4\), so it supplies no resonance through the
parameter origin. This checks the number and degrees of the local terms
required for the stated first target.

The correction uses the \emph{full analytic numerator}
\(A_S(\boldsymbol u)\) in each quotient. Replacing it by its restriction
to \(\sum_{i\in S}u_i=0\) removes the pole but changes transverse finite
coefficients. This is the same structural reason that the ambient contact
polynomials in Section~\ref{ct:section} must be differentiated before
restriction. The preserved cubic source retains these full numerators.
Its global Tornheim coordinates belong to the classical multiple-zeta
setting represented, for example, by Espinosa--Moll \cite{EspinosaMoll}.

\subsection{Results deliberately retained with their existing attribution}

The centered cubic trigamma family encountered during the audit is already
present in \emph{Quartic Spectral Reconstruction}
\cite{ProveItQuartic}, in the section containing
\src{qtm:C1}, \src{qtm:Cj}, \src{qtm:Rj}, and \src{qtm:G3}.
An independent finite summation reproduced the two exceptional
low-index corrections there. That family is not claimed as a contribution
of the present article.

Likewise, the degree bound for normalized spectral jets, the top local
polarization, and the all-difference formula already occur in
\emph{Arithmetic Spectral Jets}
\cite{ProveItArithmeticJets}, labels
\src{asj:thm:polynomial},
\src{asj:thm:top}, and
\src{asj:thm:all-differences}.
Equation~\eqref{sp:inherited-correction} is their specialized
normalization in the quadratic setting. The new ordinary density
representations should be integrated alongside this existing theorem,
not used to replace its attribution.

The unrestricted unequal-order reconstruction problem and the full
fixed-gap-six family remain open at the scope inspected here. None of the
contact or spectral identities in this article establishes their required
uniform nonvanishing. No further false theorem was established by this
targeted audit, and no claim is made that every proof in the complete
manuscript has been rechecked.

% END sections/06_audit.tex


% BEGIN sections/07_research.tex
\section{Further research questions and proposed next investigations}
\label{rq:section}

The following questions begin beyond the proved statements. They seek exact
identities, structural classifications, and finite evaluations. In
particular, a failure to recognize a numerical value in a short constant
basis is not a result about nonreducibility.

\begin{enumerate}[label=\textbf{Q\arabic*.}]
\item \textbf{A short recurrence for the full bivariate contact sequence.}
The coefficient algorithm \eqref{ct:recursion}--\eqref{ct:dk} computes every
contact, and \eqref{ct:topplus}--\eqref{ct:topminus} evaluate their highest
spectral degrees. Find a recurrence in the contact index \(m\), with
polynomial coefficients in \(p,q,a-b\), that preserves every lower-degree
Gamma correction. A satisfactory recurrence should recover the diagonal
Euler-polynomial law and exact cutoff translation without expanding all
intermediate shifted factors.

\item \textbf{Exact factorization and zero loci of the contact polynomials.}
Parity explains some factors of \(c_m,d_m\), but it does not classify all
special parameter values where a contact vanishes. Determine the complete
factorization pattern on \(p+q=-2m\), separating factors forced by exchange
symmetry from additional algebraic factors. The first target is a uniform
description of the residual factors for \(m=4,5,6\), proved from the
coefficient generator rather than inferred from numerical roots.

\item \textbf{An independent outer spectral weight.}
Adjoin \((n+c)^{-w}\) to a product of independently ordered ordinary or
alternating Hurwitz tails. Determine the moving polar hyperplanes and give
a finite-cutoff identity with a precisely specified remainder coordinate.
The present Gamma normalization chooses its outer weight through the
orders \(p,q\); it does not solve the additional independent variable
\(w\). Mellin integration of the completed endpoint law is a natural
starting point, but its boundary terms must be retained.

\item \textbf{Cyclotomic contact matrices.}
For a primitive \(d\)th-root twist with \(d>2\), replace the two-step
alternating recurrence by its \(d\)-component system. Seek a Gamma-product
normalization and an exact matrix-valued bilinear endpoint law. Determine
the resulting resonance lattice, every independent-order contact, and the
effect of shifting the cutoff by one lattice step. The case \(d=3\) is a
concrete first target; the alternating case supplies the normalization test.

\item \textbf{Three genuinely independent tails.}
For three shifts, three orders, and independent colors, classify the
smallest adjacent-product certificate whose difference removes all
unwanted residual terms. If depth-two sums are unavoidable within the
chosen certificate class, give their exact coefficients and endpoint
terms. The aim is a specified finite identity, with the size of its
function alphabet proved within that class; powers of one tail are not
sufficient evidence for the independent-data case.

\item \textbf{One ordinary full-jet integral for odd polynomial weights.}
The compact connected transform already handles arbitrary polynomial
weights, and the full even-weight jets admit the ordinary master derived
here. For an odd weight, produce a comparably economical full-jet formula
that retains all resonant Stieltjes coefficients and the inherited local
correction. A first target is \(W(x)=x\), two masses, and the third jet.
The subtraction must state its independent regulator and agree with
\eqref{sp:inherited-correction}; the known local degree theorem itself is
not the missing result.

\item \textbf{Finite special-value reductions on explicit mass loci.}
At \(a=1/2\), investigate the cubic compact integral
\eqref{da:half-cubic} and the scalar third jet at a mass satisfying
\(e^{-2\pi\sqrt\mu}=1/2\). Seek a proved reduction in a declared alphabet
of classical or multiple polylogarithms, their order derivatives, and
Gamma or zeta derivatives. Keep a finite evaluation separate from the
ordinary integral representation already established. If an obstruction
is proposed, formulate and prove it in a precise function or
difference-equation class.

\item \textbf{Complex masses and controlled branch crossings.}
Continue the density formulas to nonreal masses with explicitly chosen
contours, including conjugate pairs. Determine exact jump formulas when a
mass crosses a negative lattice value \(- (n+a)^2\). The expected jump
should be expressible through the affected finite spectral terms, but
the contour deformation and spectral normalization must be proved
compatible. This would extend the real-mass results to a useful complex
parameter geometry.

\item \textbf{Independent rational-shift projections.}
Apply finite residue-class or character projections to the all-order
digamma kernels, and compare the resulting Dirichlet-\(L\) derivatives
with the polylogarithmic density transform. Produce a complete dictionary
at conductors three and four, including the shift primitives and
coincident-mass coefficients. Specify both parity components; a
reflection identity for their sum does not determine them separately.

\item \textbf{Arithmetic reduction of selected cubic Stieltjes coefficients.}
Starting from the already-proved finite Tornheim germ, compute an explicit
small-index table containing
\(\FP\int_0^1\gamma_1^2\gamma_0'\) and its first independent
argument-derivative companions. Determine all relations obtainable from
shuffle, stuffle, symmetry, and differentiation in a declared formal
alphabet. The unresolved target is a further arithmetic reduction or an
exact formal dimension statement, not existence of the generating germ.

\item \textbf{The full reconstruction and gap-six problems.}
Retain the manuscript's unrestricted unequal-even-order reconstruction
question as a separate structural problem. A useful intermediate result
would be a bordered-Hankel identity for the entire gap-six family,
including the moving \(B_0,B_1\) terms and omitted-binomial conventions.
The known eventual nonvanishing mechanism and finite-order completions
must both be recovered. Another finite table would not resolve the
uniform question.

\item \textbf{Formal proofs and validated evaluation of the new masters.}
Formalize first the finite endpoint telescope, the terminating contact
coefficient algebra, and the polynomial subtraction identities.
Then formalize the uniform asymptotic remainder and the two Cauchy
representations with their absolute-Fubini arguments. Separately build
interval evaluators that split every logarithmic endpoint and use
analytically controlled tails. Exact polynomial checks, floating-point
diagnostics, and analytic proof have different roles and should remain
separate in the repository.
\end{enumerate}

The most direct continuation of the present work is Q6: it asks for the
remaining parity sector of an already explicit full-jet representation.
Q1 and Q4 develop the completed contact algebra structurally. Q7 and Q10
instead aim at finite arithmetic reductions, where the existing exact
integral and coefficient formulas provide reliable starting values and
normalization constraints.

% END sections/07_research.tex


% BEGIN sections/08_verification.tex
\section{Proof checks, numerical evidence, and reproducibility}
\label{vf:section}

The universal statements in this article are proved analytically:
finite telescoping identities, finite polynomial coefficient extraction,
uniform asymptotic remainders, normally convergent spectral subtractions,
and Cauchy representations with absolute-Fubini estimates. The verification
programs check finite algebra and independent numerical consequences.
They do not replace the arguments for arbitrary derivative order, contact
index, mass count, or polynomial degree.

\subsection{Exact algebra}

The contact program verifies the centered \(t_1,t_2\) coefficients, the
ambient and restricted contact tables through \(m=3\), cutoff transport
through \(m=4\), and parity, degree, and highest-homogeneous identities
through \(m=6\). It also compares the independent formula with the
predecessor's diagonal Euler-polynomial law through \(m=6\).
Together with the first-Stieltjes example program, this gives
\textbf{123 passed exact symbolic assertions}. The assertions are listed
individually in the two exact-result JSON files.

These checks operate over rational polynomial rings wherever appropriate.
An equality of expanded symbolic polynomials is exact. Substitution of a
large integer or a high-precision numerical value is not used as a
replacement for such equality. The general degree and translation theorems
are proved in Section~\ref{ct:section}, beyond the finite range checked.

\subsection{Independent decimal diagnostics}

Separate implementations check endpoints, ordinary sums, cut densities,
and arithmetic transforms. The recorded outputs include the working
precision, parameters, method, and residual, so that a small residual
is not detached from its numerical context.

\begin{center}
\small
\renewcommand{\arraystretch}{1.3}
\begin{tabular}{>{\raggedright\arraybackslash}p{0.34\textwidth}>{\raggedright\arraybackslash}p{0.34\textwidth}>{\raggedright\arraybackslash}p{0.23\textwidth}}
\toprule
Check & Independent comparison & Recorded residual scale\\
\midrule
Two \(\pi,\log2\) boundary sums &
257 direct terms plus an endpoint expansion, versus exact constants &
Below \(2.1\times10^{-27}\)\\
First-Stieltjes boundary sum &
257 direct terms plus a degree-ten endpoint expansion &
\(1.77\times10^{-28}\)\\
Stieltjes expression for \(T(z)\) &
Generalized-Stieltjes difference versus a differentiated Laplace integral &
About \(10^{-81}\) at four shifts\\
Compact cubic density, distinct and equal masses &
Direct logarithmic sums and subtracted Hurwitz tails versus digamma integrals &
Below \(1.4\times10^{-60}\)\\
Full scalar second and third jets &
Finite head plus Hurwitz-derivative tail versus the full cut integral &
\(1.64\times10^{-55}\), \(5.03\times10^{-47}\)\\
Polylogarithmic density transform &
Strict harmonic series versus quadrature at two interior points &
Below \(7\times10^{-66}\)\\
The integral \eqref{da:pi-four} &
Direct endpoint quadrature versus \(\pi^4/15\) &
\(1.71\times10^{-63}\)\\
\bottomrule
\end{tabular}
\end{center}

The contact suite records 24 numerical diagnostics at 85 working digits;
the first-Stieltjes suite records another eight at 80 digits. The compact
density tests use 60 digits, the full scalar tests 55, and the
polylogarithm tests 65. Generic endpoint-jet asymptotic checks have residuals
of the size of their first omitted terms, rather than of the nominal
working precision. The exact analytic endpoint law is tested separately.

The weighted compact tests illustrate a practical issue. Direct numerical
acceleration of a summand formed by subtracting two nearly equal large
terms initially lost several digits. The final implementation replaces that
tail with its independently convergent Hurwitz series. Its recorded
analytic bounds for the omitted power tail are below \(3\times10^{-79}\)
in the supplied cases. Those tail estimates do not enclose all rounding
and quadrature errors: the complete computation is still a diagnostic,
not an interval certificate.

The all-negative-integer extension is checked separately by comparing
its finite-subtraction continuation, an independent binomial Hurwitz
continuation, and stable quadrature of the ordinary cut integral.
At \(a=0.7\), the one-mass third derivative at \(-1\), with
\(\mu=0.8\), agrees within \(6.81\times10^{-55}\); the two-mass
second derivative at \(-1\), with \(\boldsymbol\mu=(0.8,1.4)\),
agrees within \(3.17\times10^{-52}\), at 60 working digits.
The corresponding result file records these comparisons without
conflating the analytic remainder estimate with a rigorous enclosure
of the entire floating-point computation.

\subsection{Independent analytic review}

The following sensitive steps were recomputed independently of the first
derivation: the contact translation sign; differentiation of ambient
binomial coefficients; the highest-degree hyperbolic generating functions;
the ordinary alternating subtraction; both elementary boundary constants;
the full and compact cut-density signs; the spectral scale \(2N\); the
critical odd-weight subtraction; and the Gamma normalization of the first
jet. The polynomial-weight continuation and the negative-integer cubic
example were then checked against a second derivation.

This is an internal mathematical cross-check, not a claim of external peer
review or formal certification. The finite source correction was separately
verified from the reciprocal-Gamma exponential and checked as an applicable
patch against the pinned repository.

\subsection{Using the accompanying package}

The package contains a self-contained \LaTeX{} file and its PDF, together
with modular section sources for integration. The \src{build.py} script
recreates the self-contained source and compiles it with standard
\src{pdflatex} packages. Python verification uses SymPy and mpmath; the
versions used for the recorded run appear in the result files and
\src{requirements.txt}.

Run the mathematical checks with
\begin{quote}\ttfamily
 python3 verification/run\_all.py
\end{quote}
and rebuild the article with
\begin{quote}\ttfamily
 python3 build.py
\end{quote}
The check runner writes its own machine-readable summary and preserves
separate exact and decimal records. The distributed-file manifest verifies
the byte identities of the packaged files; rebuilding a PDF may change
its metadata and hence its hash without changing the mathematical content.

The integration directory contains the single-line correction patch and
proposed status changes. The source manifest records the seven incoming
archives and the directly inspected manuscript and preserved-continuation
files. No new numerical relation is promoted to a theorem merely because
it appears in one of the diagnostic files. The unresolved questions are
listed in Section~\ref{rq:section}.

% END sections/08_verification.tex


% BEGIN references.tex
\begin{thebibliography}{99}
\raggedright

\bibitem{ProveIt}
V. Reshetnikov, \emph{ProveIt}, polylogarithm manuscript and research
continuations. Repository snapshot \pin, inspected 11 October 2026.
\href{https://github.com/VladimirReshetnikov/ProveIt/tree/e3d7ec5d3f9be259c1f91edaea4389f4e79b5b03/Analysis/Polylogarithms/docs/manuscript}{Pinned manuscript directory}.
Exact source hashes and the scope of inspection are in
\src{data/source_manifest.json}.

\bibitem{ProveItBilinear}
\emph{Bilinear Alternating Tails and Harmonic Determinants:
Rational certificates, Gamma identities, cutoff jets, and complete
order-ten reconstruction}, ProveIt research continuation, 11 October 2026.
Archive \src{ProveIt_Bilinear_Tails_and_Determinants_2026-10-11.zip}.
Especially \src{sections/02_bilinear_gamma.tex},
\src{sections/04_cutoff_jets.tex}, and the question
\emph{Every bilinear contact hyperplane} in
\src{sections/10_research_questions.tex}.
\href{https://github.com/VladimirReshetnikov/ProveIt/blob/e3d7ec5d3f9be259c1f91edaea4389f4e79b5b03/docs/incoming/ProveIt_Bilinear_Tails_and_Determinants_2026-10-11.zip}{Pinned archive}.

\bibitem{ProveItTailKernels}
\emph{Bilinear Tail Identities, Polylogarithmic Resonances, and Stieltjes
Boundary Products: Complete moment reconstruction at orders ten and twelve},
ProveIt research continuation, 11 October 2026.
Archive \src{ProveIt_Tail_Kernels_and_Boundary_Products_2026-10-11.zip}.
Especially research questions R7 and R9 in
\src{sections/08_audit_and_research.tex}.
\href{https://github.com/VladimirReshetnikov/ProveIt/blob/e3d7ec5d3f9be259c1f91edaea4389f4e79b5b03/docs/incoming/ProveIt_Tail_Kernels_and_Boundary_Products_2026-10-11.zip}{Pinned archive}.

\bibitem{ProveItHigherQuadratic}
\emph{Higher Quadratic Spectral Jets: Harmonic--Stieltjes factorization
defects, polylogarithmic second derivatives, and exact mass collisions},
ProveIt research continuation, 11 October 2026.
Archive \src{ProveIt_Higher_Quadratic_Jets_2026-10-11.zip}.
Especially \src{sections/10-research.tex}, items 1, 3, and 7.
\href{https://github.com/VladimirReshetnikov/ProveIt/blob/e3d7ec5d3f9be259c1f91edaea4389f4e79b5b03/docs/incoming/ProveIt_Higher_Quadratic_Jets_2026-10-11.zip}{Pinned archive}.

\bibitem{ProveItHigherDefects}
\emph{Beta Polarization of Higher Spectral Defects: All-order Stieltjes
counterterms, Gamma integrals, polylogarithmic second jets, and confluent
quadratic masses}, ProveIt research continuation, 11 October 2026.
Archive \src{ProveIt_Higher_Spectral_Defects_2026-10-11.zip}.
Especially the second-jet many-mass formulas, the confluent expansion, and
further research questions 2--4.
\href{https://github.com/VladimirReshetnikov/ProveIt/blob/e3d7ec5d3f9be259c1f91edaea4389f4e79b5b03/docs/incoming/ProveIt_Higher_Spectral_Defects_2026-10-11.zip}{Pinned archive}.

\bibitem{ProveItQuartic}
\emph{Harmonic-Tail Reconstruction, Quartic Trigamma Moments, and Spectral
Stieltjes Calculus}, ProveIt research continuation, 11 October 2026.
Archive \src{ProveIt_Quartic_Spectral_Reconstruction_2026-10-11.zip}.
Especially \src{sections/04_quartic.tex} for the retained cubic family.
\href{https://github.com/VladimirReshetnikov/ProveIt/blob/e3d7ec5d3f9be259c1f91edaea4389f4e79b5b03/docs/incoming/ProveIt_Quartic_Spectral_Reconstruction_2026-10-11.zip}{Pinned archive}.

\bibitem{ProveItArithmeticJets}
\emph{Arithmetic Spectral Jets: Stieltjes residues, Gamma reciprocity,
and polylogarithmic untwisting}, preserved ProveIt continuation.
\src{stieltjes-harmonic-resonance/continuations/arithmetic-spectral-jets}.
Especially \src{sections/03-local.tex}, theorems
\src{asj:thm:polynomial}, \src{asj:thm:top}, and
\src{asj:thm:all-differences}.
\href{https://github.com/VladimirReshetnikov/ProveIt/blob/e3d7ec5d3f9be259c1f91edaea4389f4e79b5b03/Analysis/Polylogarithms/docs/reports/stieltjes-harmonic-resonance/continuations/arithmetic-spectral-jets/sections/03-local.tex}{Pinned source}.

\bibitem{ProveItCubic}
\emph{Coincident Stieltjes Products and Directional Zeta Identities},
preserved ProveIt continuation.
\src{stieltjes-correlation-calculus/continuations/coincident-directional-zeta}.
Especially \src{sections/02_cubic.tex}, labels
\src{cub:general-completion}, \src{cub:all-moments}, and \src{cub:global}.
\href{https://github.com/VladimirReshetnikov/ProveIt/blob/e3d7ec5d3f9be259c1f91edaea4389f4e79b5b03/Analysis/Polylogarithms/docs/reports/stieltjes-correlation-calculus/continuations/coincident-directional-zeta/sections/02_cubic.tex}{Pinned source}.

\bibitem{DLMF}
NIST Digital Library of Mathematical Functions.
\S5.11, \emph{Asymptotic expansions};
\S5.15, \emph{Polygamma functions};
\S25.11, \emph{Hurwitz zeta function};
\S25.12, \emph{Polylogarithms}.
\url{https://dlmf.nist.gov/5.11},
\url{https://dlmf.nist.gov/5.15},
\url{https://dlmf.nist.gov/25.11},
\url{https://dlmf.nist.gov/25.12}.
Consulted 11 October 2026.

\bibitem{HuKim}
S. Hu and M.-S. Kim,
\emph{Asymptotic expansions for the alternating Hurwitz zeta function
and its derivatives}, arXiv:2103.15528.
\url{https://arxiv.org/abs/2103.15528}.
Used for the classical expansion and derivative setting, not as a source
of the independent contact polynomials established here.

\bibitem{WangHuKim}
W. Wang, S. Hu, and M.-S. Kim,
\emph{On gamma functions with respect to the alternating Hurwitz zeta
functions}, arXiv:2405.02854.
\url{https://arxiv.org/abs/2405.02854}.
The modified Gamma quotient is classical background for \(h(z)\).

\bibitem{Hoffman}
M. E. Hoffman,
\emph{Sums of products of Riemann zeta tails},
Mediterranean Journal of Mathematics \textbf{13} (2016), 2771--2781.
arXiv:1602.03399.
\url{https://arxiv.org/abs/1602.03399}.

\bibitem{PaychaScott}
S. Paycha and S. Scott,
\emph{A Laurent expansion for regularised integrals of holomorphic symbols},
Geometric and Functional Analysis \textbf{17} (2007), 491--536.
arXiv:math/0506211.
\url{https://arxiv.org/abs/math/0506211}.
General background for local terms in regularized traces.

\bibitem{EspinosaMoll}
O. Espinosa and V. H. Moll,
\emph{The evaluation of Tornheim double sums. Part 1},
Journal of Number Theory \textbf{116} (2006), 200--229.
arXiv:math/0505647.
\url{https://arxiv.org/abs/math/0505647}.
\end{thebibliography}

% END references.tex

\end{document}
